\documentclass[a4paper,11pt]{article}
        \usepackage[top=15mm,bottom=18mm,left=16mm,right=16mm]{geometry}
        \usepackage{fontspec}
        \usepackage{xeCJK}
        \setmainfont{Times New Roman}
        \setCJKmainfont{Yu Gothic}
        \setmonofont{Consolas}
        \setCJKmonofont{Yu Gothic}
        \usepackage{booktabs}
        \usepackage{longtable}
        \usepackage{tabularx}
        \usepackage{array}
        \usepackage{enumitem}
        \usepackage{hyperref}
        \usepackage{listings}
        \usepackage{xcolor}
        \usepackage{amsmath}
        \usepackage{amssymb}
        \hypersetup{
          colorlinks=true,
          linkcolor=blue,
          urlcolor=blue,
          pdftitle={dashboard + HTTP API},
          pdfauthor={Codex}
        }
        \lstset{
          basicstyle=\ttfamily\small,
          breaklines=true,
          columns=fullflexible,
          keepspaces=true,
          frame=single,
          framerule=0.2pt,
          rulecolor=\color{black},
          showstringspaces=false
        }
        \setlength{\parskip}{0.42em}
        \setlength{\parindent}{1em}
        \setlength{\emergencystretch}{3em}
        \renewcommand{\arraystretch}{1.15}
        \title{30. dashboard + HTTP API}
        \author{solar\_ws0129-main}
        \date{2026-07-29}
        \begin{document}
        \maketitle
        \tableofcontents
        \clearpage

        \section{対象}
        \begin{tabularx}{\linewidth}{>{\raggedright\arraybackslash}p{0.18\linewidth} X}
        \toprule
        項目 & 内容 \\
        \midrule
        ファイル & \path|mpc_solarcar/dashboard_node.py| \\
        ソースSHA-256 & \path|26286dcb019288a13c2e17be4574c311b8bc2bb5e0268544165d6519a5cba443| \\
        種別 & Python \\
        区分 & runtime node \\
        役割 & planner と vehicle の現在値を ROS から受け、HTTP サーバと dashboard frontend へ渡す。 \\
        起動文脈 & 可視化の中心ノード。 \\
        \bottomrule
        \end{tabularx}

        \section{このファイルを読む理由}
        planner と vehicle の現在値を ROS から受け、HTTP サーバと dashboard frontend へ渡す。

        \begin{itemize}[leftmargin=1.6em]
        \item /api/state と /metrics を持つ。
\item ROS topic を直接 browser へ出すのではなく、Node 内 state に集約する。
        \end{itemize}

        \section{呼び出し関係}
        \subsection{どこから呼ばれるか}
        \begin{itemize}[leftmargin=1.6em]
        \item \path|launch/solarcar_sim.launch.py|
\item \path|mpc_solarcar/live_launch.py|
\item \path|launch/solar_measurement.launch.py|
        \end{itemize}

        \subsection{次に読むと理解が進むファイル}
        \begin{itemize}[leftmargin=1.6em]
        \item 特記事項なし。
        \end{itemize}

        \subsection{同一リポジトリ内の主要依存}
        \begin{itemize}[leftmargin=1.6em]
        \item 同一リポジトリ内の明示 import は少ない。
        \end{itemize}

        \section{ソース冒頭の把握}
        モジュール docstring または先頭意図の要約:

        モジュール docstring は明示されていない。

        \subsection{代表コード断片}
        \begin{lstlisting}[language=Python]
from ament_index_python.packages import get_package_share_directory


class DashboardHandler(SimpleHTTPRequestHandler):
    def __init__(self, *args, node=None, directory=None, **kwargs):
        self._node = node
        super().__init__(*args, directory=directory, **kwargs)

    def do_GET(self):
        if self.path.startswith('/api/state'):
            self._send_json(self._node.get_state())
            return
        if self.path.startswith('/metrics'):
            self._send_metrics(self._node.get_prometheus_metrics())
            return
        super().do_GET()

    def log_message(self, format, *args):
        # quiet
        return

    def _send_json(self, data):
\end{lstlisting}


        \section{主要構造の読み取り}
        主要クラスは DashboardHandler, DashboardNode。 主要関数は do\_GET, log\_message, destroy\_node, get\_state, get\_prometheus\_metrics, main。 ROS パラメータ宣言は 5 件。 ROS I/O は publisher 0 件、subscription 21 件。

        \subsection{主要クラス・関数}
            \begin{longtable}{p{0.07\linewidth} p{0.16\linewidth} p{0.25\linewidth} p{0.40\linewidth}}
            \toprule
            行 & 種別 & 名前 & 補足 \\
            \midrule
            19 & class & \path|DashboardHandler| & bases: SimpleHTTPRequestHandler \\
55 & class & \path|DashboardNode| & bases: Node \\
20 & function & \path|__init__| & args: self \\
24 & function & \path|do_GET| & args: self \\
33 & function & \path|log_message| & args: self, format \\
37 & function & \path|_send_json| & args: self, data \\
45 & function & \path|_send_metrics| & args: self, payload \\
56 & function & \path|__init__| & args: self \\
156 & function & \path|destroy_node| & args: self \\
165 & function & \path|get_state| & args: self \\
169 & function & \path|get_prometheus_metrics| & args: self \\
236 & function & \path|_update| & args: self, key, value \\
241 & function & \path|_on_speed_cmd| & args: self, msg \\
244 & function & \path|_on_upper_speed_cmd| & args: self, msg \\
247 & function & \path|_on_speed_meas| & args: self, msg \\
250 & function & \path|_on_throttle_cmd| & args: self, msg \\
253 & function & \path|_on_drive_mode| & args: self, msg \\
256 & function & \path|_on_soc| & args: self, msg \\
264 & function & \path|_set_estimated_soc| & args: self, value \\
273 & function & \path|_on_tb| & args: self, msg \\
276 & function & \path|_on_ibatt| & args: self, msg \\
279 & function & \path|_on_vbatt| & args: self, msg \\
282 & function & \path|_on_s_km| & args: self, msg \\
285 & function & \path|_on_status| & args: self, msg \\
301 & function & \path|_on_env| & args: self, msg \\
313 & function & \path|_on_metrics| & args: self, msg \\
333 & function & \path|_on_upper_plan| & args: self, msg \\
344 & function & \path|_on_lower_plan| & args: self, msg \\
355 & function & \path|_on_sys_state| & args: self, msg \\
358 & function & \path|_on_sys_diag| & args: self, msg \\
361 & function & \path|_on_mpc_state| & args: self, msg \\
364 & function & \path|_on_health| & args: self, msg \\
367 & function & \path|_on_gps| & args: self, msg \\
373 & function & \path|_on_path| & args: self, msg \\
379 & function & \path|_init_dummy| & args: self, path, rate\_hz \\
395 & function & \path|_tick_dummy| & args: self \\
409 & function & \path|main| &  \\
            \bottomrule
            \end{longtable}

        \subsection{ファイルを上から読んだときの定義順}
            \begin{longtable}{p{0.07\linewidth} p{0.84\linewidth}}
            \toprule
            行 & 内容 \\
            \midrule
            19 & クラス DashboardHandler を定義する。 \\
55 & クラス DashboardNode を定義する。 \\
409 & 関数 main を定義する。 \\
417 & 条件 \_\_name\_\_ == '\_\_main\_\_' を判定し、真なら内部処理を行う。 \\
            \bottomrule
            \end{longtable}

        \subsection{import 群の詳細}
            \begin{longtable}{p{0.07\linewidth} p{0.33\linewidth} p{0.50\linewidth}}
            \toprule
            行 & import 文 & このファイルで読む理由 \\
            \midrule
            1 & \path|import json| & manifest、checkpoint、UDP payload をやり取りするため。 このファイル内での主な使用位置は L38。 \\
2 & \path|import math| & Python標準のスカラー数値関数と定数を使うため。実際に参照する名前と行は使用位置欄で確認する。 このファイル内での主な使用位置は L212。 \\
3 & \path|import os| & パス、環境変数、プロセス外部状態を扱うため。 このファイル内での主な使用位置は L68, L70, L71。 \\
4 & \path|import threading| & threading モジュールを利用するため。 このファイル内での主な使用位置は L73, L150。 \\
5 & \path|import time| & wall/monotonic time に基づく周期制御や freshness 判定を行うため。 このファイル内での主な使用位置は L75, L215, L239, L261, L262, L268, L299, L311, ...。 \\
6 & \path|from functools import partial| & functools から partial を読み込み、このファイルの処理を組み立てるため。 このファイル内での主な使用位置は L147。 \\
7 & \path|from http.server import SimpleHTTPRequestHandler, ThreadingHTTPServer| & http.server から SimpleHTTPRequestHandler, ThreadingHTTPServer を読み込み、このファイルの処理を組み立てるため。 このファイル内での主な使用位置は L19, L149。 \\
9 & \path|import rclpy| & ROS 2 Python ノードとして起動・spin するため。 このファイル内での主な使用位置は L410, L412, L414。 \\
10 & \path|from rclpy.node import Node| & ROS 2 ノード本体の基底クラスとして使うため。 このファイル内での主な使用位置は L55。 \\
12 & \path|from std_msgs.msg import Float32, Float32MultiArray, String| & Float32、Bool、Stringなどの標準ROS messageで値をpublishまたはsubscribeするため。 このファイル内での主な使用位置は L117, L118, L119, L120, L121, L122, L123, L124, ...。 \\
13 & \path|from sensor_msgs.msg import NavSatFix| & GPS 緯度経度高度を ROS topic でやり取りするため。 このファイル内での主な使用位置は L136, L367。 \\
14 & \path|from nav_msgs.msg import Path| & 将来軌跡を dashboard や logger へ出すため。 このファイル内での主な使用位置は L137, L373。 \\
16 & \path|from ament_index_python.packages import get_package_share_directory| & ament\_index\_python.packages から get\_package\_share\_directory を読み込み、このファイルの処理を組み立てるため。 このファイル内での主な使用位置は L67。 \\
381 & \path|import pandas as pd| & CSV や時刻列の読込・整形・再サンプリングに使うため。 このファイル内での主な使用位置は L382。 \\
            \bottomrule
            \end{longtable}

        \section{このファイルを読む前に必要な基礎知識}
        次の章は、構文やROS用語を既知と仮定しないための説明である。

        \subsection{プログラム、プロセス、メモリ、オブジェクトを区別する}
ソースファイルはディスク上の文字列であり、それ自体は走っていない。Python実行可能プログラムがソースを読み、OSがその実行を一つのプロセスとして管理する。プロセスには仮想メモリ、開いているファイル、スレッド、終了コードなどが対応する。


\begin{lstlisting}
ソースファイル -> Pythonインタプリタが読み込む -> OS上のプロセス -> Pythonオブジェクトをメモリに生成 -> 関数やcallbackを実行
\end{lstlisting}


「メモリ上に生成する」とは、実行中プロセスが使う記憶領域に、その値の型、属性、参照関係を表す実体を用意することである。変数は実体そのものというより、そのオブジェクトを指す名前として理解するとPythonの代入が読みやすい。


\begin{lstlisting}[language=python]
node = MPCNode()
alias = node
# nodeとaliasは同じオブジェクトを参照する。
\end{lstlisting}


一つのプロセスには複数スレッドを持たせられる。スレッドは同じプロセスのメモリを共有するため、受信callbackと最適化callbackが同じself.zを書き換える場合は実行順と排他を検討する必要がある。

\paragraph{根拠資料}
\begin{itemize}[leftmargin=1.6em]
\item \href{https://docs.python.org/3/tutorial/classes.html}{Python公式チュートリアル: Classes}
\end{itemize}

\subsection{名前、代入、型、参照}
Pythonの代入文は、右辺を先に評価し、その結果のオブジェクトへ左辺の名前を結び付ける。`x = f()`では、まずfを呼び、その戻り値が得られてからxが更新される。


`float(...)`、`int(...)`、`bool(...)`は型名を呼び出して値を変換する。外部CSV、YAML、ROS parameterから来た値は型が期待どおりとは限らないため、このリポジトリでは明示変換が多い。


`None`は値が無いことを示す単一のオブジェクト、`math.nan`は浮動小数点値ではあるが有効な数値ではないことを示す。両者は用途が異なるため、`is None`と`math.isfinite`を使い分ける。

\paragraph{根拠資料}
\begin{itemize}[leftmargin=1.6em]
\item \href{https://docs.python.org/3/tutorial/classes.html}{Python公式チュートリアル: Classes}
\end{itemize}

\subsection{関数、引数、戻り値、スコープ}
`def`は関数本体をその場で実行する文ではない。関数オブジェクトを作り、その名前を現在の名前空間へ登録する。`f`は関数そのもの、`f()`はその関数を呼んだ結果である。


仮引数は関数定義側の受け取り名、実引数は呼出側から渡す値である。位置引数、キーワード引数、既定値付き引数、`*args`、`**kwargs`は渡し方が異なる。


\begin{lstlisting}[language=python]
def clip_speed(value, minimum=0.0, maximum=110.0):
    return min(maximum, max(minimum, value))

v = clip_speed(120.0, maximum=90.0)
\end{lstlisting}


関数内で作った通常の名前はローカルスコープに属する。メソッドが`self.z`を更新すればオブジェクトの状態が残るが、単に`z`を更新しただけならその呼出しのローカル値である。

\paragraph{根拠資料}
\begin{itemize}[leftmargin=1.6em]
\item \href{https://docs.python.org/3/tutorial/classes.html}{Python公式チュートリアル: Classes}
\end{itemize}

\subsection{module、package、import、console\_scripts}
Pythonファイルはmoduleとして読み込める。複数moduleをディレクトリにまとめたものがpackageである。`from .model import SolarCarModel`の先頭の点は、現在と同じpackage内のmodel moduleを指す相対importである。


import時にはファイルのトップレベル文が上から一度実行される。`def`や`class`は関数・クラスを登録するが、その本体の通常処理は呼び出すまで走らない。


setuptoolsの`console\_scripts`は、端末で使う実行可能名と`package.module:function`を対応付けるインストール時メタデータである。生成される実行用ラッパーはmoduleをimportし、指定関数を呼び、戻り値を終了コードとして扱う小さな入口である。


\begin{lstlisting}
ROS launchのexecutable名 -> install済みconsole script -> mpc_solarcar.mpc_nodeをimport -> main()を呼ぶ
\end{lstlisting}

\paragraph{根拠資料}
\begin{itemize}[leftmargin=1.6em]
\item \href{https://setuptools.pypa.io/en/latest/userguide/entry_point.html}{setuptools公式: Entry Points / Console Scripts}
\item \href{https://docs.python.org/3/tutorial/classes.html}{Python公式チュートリアル: Classes}
\end{itemize}

\subsection{例外、try/finally、with、資源解放}
例外は通常の戻り値とは別経路で異常を呼出元へ伝える。`try/except`は想定した異常を処理し、`finally`は成功・失敗にかかわらず後始末を行う。


`with open(...) as f:`はcontext managerを使い、ブロックを出るとファイルを閉じる。CSVやログの破損を避けるため、開いた資源の所有者と閉じる場所を明確にする。


`except Exception: pass`はノードを止めない利点がある一方、入力異常を隠して原因追跡を難しくする。安全に関係する値では、少なくとも頻度制限付き警告、異常カウンタ、fallback状態のpublishを検討する。



\subsection{class、独自クラス、インスタンス、self、継承}
クラスは、保持するデータと、そのデータを扱う関数を一つの型としてまとめる設計単位である。「独自クラス」とはPythonやROS 2が最初から用意した型ではなく、このプロジェクトが目的に合わせて定義した新しい型を指す。


`class MPCNode(Node)`は、rclpyのNodeを基底クラスとして継承し、Nodeが持つpublisher、subscription、timer、parameterなどの機能にMPC固有機能を追加する。丸括弧内は関数引数ではなく基底クラス指定である。


`MPCNode()`はクラスオブジェクトを呼び出して新しいインスタンスを作る。Pythonは新しい実体を用意し、その実体を第1引数selfとして`\_\_init\_\_`へ渡す。`self`は予約語ではなく慣習的な名前だが、この慣習を守ることでコードの意味が共有される。


\begin{lstlisting}[language=python]
class Counter:
    def __init__(self):
        self.value = 0

    def increment(self):
        self.value += 1

counter = Counter()
counter.increment()
\end{lstlisting}


`super().\_\_init\_\_(...)`は継承元の初期化処理を呼ぶ。MPCNodeではこれを省くとROSノードとして必要な内部実体が作られず、`create\_publisher`などを正しく使えない。

\paragraph{根拠資料}
\begin{itemize}[leftmargin=1.6em]
\item \href{https://docs.python.org/3/tutorial/classes.html}{Python公式チュートリアル: Classes}
\item \href{https://docs.ros.org/en/humble/Tutorials/Beginner-CLI-Tools/Understanding-ROS2-Nodes/Understanding-ROS2-Nodes.html}{ROS 2 Humble公式: Understanding nodes}
\end{itemize}

\subsection{先頭アンダースコア、\_\_init\_\_、名前付けの作法}
`self.\_step\_solar`の先頭1個のアンダースコアは「クラス内部で使う実装詳細」という慣習であり、アクセスを強制禁止する機構ではない。外部コードから呼べるが、公開APIとして依存しない意思表示である。


`\_\_init\_\_`の前後2個のアンダースコアはPythonが定めた特殊メソッド名である。クラス生成時に自動で呼ばれる。任意の名前に前後2個を付けて独自仕様を作ることは避ける。


`\_`で始まるローカル変数は、未使用または外部へ見せない意図を示す場合がある。例えば`\_base\_csv`は戻り値として受け取る必要はあるが、その後の処理では使わないことを読者へ示す。

\paragraph{根拠資料}
\begin{itemize}[leftmargin=1.6em]
\item \href{https://docs.python.org/3/tutorial/classes.html}{Python公式チュートリアル: Classes}
\end{itemize}

\subsection{list、dict、deque、NumPy配列、pandas DataFrame}
listは順序付き可変列、tupleは順序付きで通常変更しない列、dictはキーから値を引く対応表、dequeは両端追加・削除が効率的なキューである。どの構造を選ぶかはアクセス方法と更新方法で決まる。


NumPy配列は同種数値を連続的に扱い、要素ごとの演算、clip、補間、線形代数を簡潔に書く。shapeは各次元の要素数であり、速度系列なら通常`(N,)`、候補集団なら`(population, N)`となる。


pandas DataFrameは列名を持つ表である。CSV読込後は列型、欠損、単位、timezone、並び順、重複を明示的に処理しなければ、数値計算が動いても意味が誤る。



\subsection{rclpy、rcl、rmw、DDS/RTPSの層}
rclpyはPython利用者向けROS 2 client libraryである。利用者のNode、publisher、subscriptionなどをPython APIとして提供し、その下で共通C層のrclを利用する。


\begin{lstlisting}
本プロジェクトのPythonコード -> rclpy -> rcl -> rmw -> DDS/RTPS実装 -> ネットワークまたは同一PC内通信
\end{lstlisting}


rclは言語に依存しない共通ROS機能を提供するC API、rmwはROS 2と具体的middleware実装の境界である。DDS/RTPS側が探索、serialize、publish/subscribe、request/replyなどを担う。executorの実行モデルはrclだけで完結せずclient library側にも実装される。

\paragraph{根拠資料}
\begin{itemize}[leftmargin=1.6em]
\item \href{https://docs.ros.org/en/humble/Concepts/About-Internal-Interfaces.html}{ROS 2 Humble公式: Internal ROS 2 interfaces}
\end{itemize}

\subsection{ROS graph、Node、topic、service、action、parameter}
ROS graphは、実行中Nodeと、それらが持つpublisher、subscription、service、actionなどの接続関係である。Pythonクラス、プロセス、ROS graph上のNode名は関連するが同一物ではない。


topicは継続データ向けの非同期一方向publish/subscribe、serviceは短いrequest/response、actionは時間のかかる目標へfeedback、cancel、resultを持たせる。parameterはNodeごとの設定値である。


publisherが送った時点とsubscription callbackが実行される時点は同じとは限らない。通信遅延、QoS queue、executor待ちを挟むため、センサ値にはtimestampとfreshness判定が必要になる。

\paragraph{根拠資料}
\begin{itemize}[leftmargin=1.6em]
\item \href{https://docs.ros.org/en/humble/Tutorials/Beginner-CLI-Tools/Understanding-ROS2-Nodes/Understanding-ROS2-Nodes.html}{ROS 2 Humble公式: Understanding nodes}
\item \href{https://docs.ros.org/en/humble/Concepts/Basic/Interfaces-Topics-Services-Actions.html}{ROS 2 Humble公式: Topics, Services, Actions}
\item \href{https://docs.ros.org/en/humble/Concepts/Basic/About-Parameters.html}{ROS 2 Humble公式: Parameters}
\end{itemize}

\subsection{callback、timer、Executor、spin、Callback Group}
callbackは、メッセージ受信、timer満了、service requestなどのイベントが成立した後でExecutorから呼ばれる関数である。登録時に`self.\_on\_speed`と括弧なしで渡すのは、今実行せず後で呼ぶ関数オブジェクトを渡すためである。


Executorは実行可能になったcallbackを見つけ、Callback Groupの条件を確認して実行する。`spin()`は終了要求まで待機・dispatchを続けるため、通常運転中はmainの次の行へ戻らない。


MutuallyExclusiveCallbackGroupは同じgroup内のcallbackを同時実行させない。別group間はMultiThreadedExecutorで並行実行し得る。groupを分けただけでは共有属性への完全な排他にはならないため、複数groupが同じ状態を読む・書く場合は設計確認が必要である。


\begin{lstlisting}
DDS受信またはtimer満了 -> wait setでready -> Executorが選択 -> Callback Groupが許可 -> worker threadがcallback実行 -> 終了後Executorへ戻る
\end{lstlisting}

\paragraph{根拠資料}
\begin{itemize}[leftmargin=1.6em]
\item \href{https://docs.ros.org/en/rolling/Concepts/Intermediate/About-Executors.html}{ROS 2公式: Executors}
\item \href{https://docs.ros.org/en/humble/Concepts/About-Internal-Interfaces.html}{ROS 2 Humble公式: Internal ROS 2 interfaces}
\end{itemize}

\subsection{rqt\_graph、ros2 CLI、rosbag2をいつ使うか}
rqt\_graphは接続関係を見る道具であり、数値の正しさや更新周期までは保証しない。起動直後、topic名変更後、publisherが複数存在する疑いがある時に使う。


\begin{lstlisting}[language=bash]
ros2 node list
ros2 node info /mpc_node
ros2 topic list -t
ros2 topic info -v /planner/speed_cmd
ros2 topic hz /planner/speed_cmd
ros2 topic echo /planner/status
rqt_graph
\end{lstlisting}


rosbag2はtopicメッセージを時系列のまま記録・再生する。通信不具合、freshness、再計画trigger、実車とSILSの差を再現可能にするため、本番前試験では制御入力だけでなく原因となる全telemetry、status、parameter情報を記録する。


\begin{lstlisting}[language=bash]
ros2 bag record -o outputs/bags/preflight \
  /vehicle/s_km /vehicle/speed_kmh /vehicle/batt_soc \
  /vehicle/batt_temp_c /vehicle/batt_current_a /vehicle/batt_voltage_v \
  /planner/upper_speed_cmd /planner/speed_cmd /planner/status

ros2 bag info outputs/bags/preflight
ros2 bag play outputs/bags/preflight --clock
\end{lstlisting}


bag再生時はQoS互換性、simulation time、外部publisherとの二重入力に注意する。実車Nodeを同時に動かす場合はnamespaceまたはremappingで入力源を明確に分離する。

\paragraph{根拠資料}
\begin{itemize}[leftmargin=1.6em]
\item \href{https://docs.ros.org/en/ros2_packages/humble/api/rqt_graph/index.html}{ROS 2 Humble公式: rqt\_graph}
\item \href{https://docs.ros.org/en/humble/Tutorials/Beginner-CLI-Tools/Recording-And-Playing-Back-Data/Recording-And-Playing-Back-Data.html}{ROS 2 Humble公式: Recording and playing back data}
\item \href{https://docs.ros.org/en/humble/Tutorials/Beginner-CLI-Tools/Understanding-ROS2-Nodes/Understanding-ROS2-Nodes.html}{ROS 2 Humble公式: Understanding nodes}
\end{itemize}

\subsection{天候、route、補間、時刻、単位}
予報は時刻だけの系列か、時刻とroute距離の2次元gridになり得る。現在時刻tと距離sがgrid点の間なら、周囲の値から補間して日射、温度、風を求める。


UTC、現地timezone、naive datetimeを混ぜると数時間のずれが発生する。入力でtimezoneを確定し、内部比較はUTC、表示は現地時刻という役割分担が安全である。


route距離のkm、速度のkm/hとm/s、時間の秒、energyのWhとJを跨ぐため、変換係数3.6、3600、1000の意味を各式で確認する。



\subsection{制御、状態、入力、モデル予測制御MPC}
制御対象の内部を表す状態をx、操作入力をu、外乱・予報をwとすると、離散モデルは`x[k+1] = f(x[k], u[k], w[k])`と書ける。ソーラーカーではSoC、電池温度、距離、速度などが状態候補、速度目標や駆動トルクが入力候補になる。


\[
\min_{u_0,\ldots,u_{N-1}} \sum_{k=0}^{N-1}\ell(x_k,u_k,w_k)+V_f(x_N)
\]

MPCは現在状態からNステップ先まで予測し、目的関数と制約を満たす入力系列を求める。ただし実際に適用するのは通常先頭入力だけで、次回は新しい実測状態から再び解く。これがreceding horizonである。


予測モデル、目的関数、制約、ホライズン、solver、初期値のどれかが変わると答えも変わる。「MPCを使う」だけでは仕様は決まらず、これらを単位付きで追う必要がある。

\paragraph{根拠資料}
\begin{itemize}[leftmargin=1.6em]
\item \href{https://docs.scipy.org/doc/scipy/reference/generated/scipy.optimize.minimize.html}{SciPy公式: scipy.optimize.minimize}
\end{itemize}



        \section{関数・クラスを上から順に解説}
        \subsection{L19 クラス DashboardHandler}
            \begin{itemize}[leftmargin=1.6em]
              \item 行範囲: L19--L52
              \item 所属: module直下
              \item 定義: \path|DashboardHandler(bases=SimpleHTTPRequestHandler)|
              \item docstring: 明示されていない。
              \item このブロックが直接呼ぶ主な関数・メソッド: 特に明示されていない。
              \item 戻り値の要点: 戻り値が無いか、return 文が明示されていない。
              \item 主なローカル代入名: 特になし。
              \item 読み取る主なインスタンス属性: 特になし。
              \item 更新する主なインスタンス属性: 特になし。
              \item 明示的に送出する例外: 明示的raiseはない。
              \item 制御構造の規模: 条件分岐 0、ループ 0、try 0
            \end{itemize}
            \paragraph{この定義を読むためのPython構文}
            \begin{itemize}[leftmargin=1.6em]
            \item class文は新しい型を定義する。丸括弧内は継承する基底クラスである。
            \end{itemize}
            \paragraph{上から順の処理}
            \begin{enumerate}[leftmargin=1.8em]
            \item 関数 \_\_init\_\_ を定義する。
\item 関数 do\_GET を定義する。
\item 関数 log\_message を定義する。
\item 関数 \_send\_json を定義する。
\item 関数 \_send\_metrics を定義する。
            \end{enumerate}
            \paragraph{代表コード断片}
            \begin{lstlisting}[language=Python]
class DashboardHandler(SimpleHTTPRequestHandler):
    def __init__(self, *args, node=None, directory=None, **kwargs):
        self._node = node
        super().__init__(*args, directory=directory, **kwargs)

    def do_GET(self):
        if self.path.startswith('/api/state'):
            self._send_json(self._node.get_state())
            return
        if self.path.startswith('/metrics'):
            self._send_metrics(self._node.get_prometheus_metrics())
            return
        super().do_GET()

    def log_message(self, format, *args):
        # quiet
        return

    def _send_json(self, data):
        payload = json.dumps(data).encode('utf-8')
        self.send_response(200)
        self.send_header('Content-Type', 'application/json')
        self.send_header('Content-Length', str(len(payload)))
        self.end_headers()
        self.wfile.write(payload)

    def _send_metrics(self, payload):
        body = payload.encode('utf-8')
        self.send_response(200)
        self.send_header('Content-Type', 'text/plain; version=0.0.4; charset=utf-8')
        self.send_header('Cache-Control', 'no-store')
        self.send_header('Content-Length', str(len(body)))
        self.end_headers()
        self.wfile.write(body)
\end{lstlisting}

\subsection{L20 関数 DashboardHandler.\_\_init\_\_}
            \begin{itemize}[leftmargin=1.6em]
              \item 行範囲: L20--L22
              \item 所属: \path|DashboardHandler|
              \item 定義: \path|__init__(self, *args, node = None, directory = None, **kwargs)|
              \item docstring: 明示されていない。
              \item このブロックが直接呼ぶ主な関数・メソッド: \_\_init\_\_, super
              \item 戻り値の要点: 戻り値が無いか、return 文が明示されていない。
              \item 主なローカル代入名: 特になし。
              \item 読み取る主なインスタンス属性: 特になし。
              \item 更新する主なインスタンス属性: self.\_node
              \item 明示的に送出する例外: 明示的raiseはない。
              \item 制御構造の規模: 条件分岐 0、ループ 0、try 0
            \end{itemize}
            \paragraph{この定義を読むためのPython構文}
            \begin{itemize}[leftmargin=1.6em]
            \item 第1引数selfは、このメソッドを呼んだインスタンス自身を受け取る慣習名である。
            \end{itemize}
            \paragraph{上から順の処理}
            \begin{enumerate}[leftmargin=1.8em]
            \item self.\_node に node の結果を代入する。
\item super().\_\_init\_\_(...) を実行する。
            \end{enumerate}
            \paragraph{代表コード断片}
            \begin{lstlisting}[language=Python]
    def __init__(self, *args, node=None, directory=None, **kwargs):
        self._node = node
        super().__init__(*args, directory=directory, **kwargs)
\end{lstlisting}

\subsection{L24 関数 DashboardHandler.do\_GET}
            \begin{itemize}[leftmargin=1.6em]
              \item 行範囲: L24--L31
              \item 所属: \path|DashboardHandler|
              \item 定義: \path|do_GET(self)|
              \item docstring: 明示されていない。
              \item このブロックが直接呼ぶ主な関数・メソッド: \_send\_json, \_send\_metrics, do\_GET, get\_prometheus\_metrics, get\_state, startswith, super
              \item 戻り値の要点: 戻り値が無いか、return 文が明示されていない。
              \item 主なローカル代入名: 特になし。
              \item 読み取る主なインスタンス属性: self.\_node, self.\_send\_json, self.\_send\_metrics, self.path
              \item 更新する主なインスタンス属性: 特になし。
              \item 明示的に送出する例外: 明示的raiseはない。
              \item 制御構造の規模: 条件分岐 2、ループ 0、try 0
            \end{itemize}
            \paragraph{この定義を読むためのPython構文}
            \begin{itemize}[leftmargin=1.6em]
            \item 第1引数selfは、このメソッドを呼んだインスタンス自身を受け取る慣習名である。
            \end{itemize}
            \paragraph{上から順の処理}
            \begin{enumerate}[leftmargin=1.8em]
            \item 条件 self.path.startswith('/api/state') を判定し、真なら内部処理を行う。
\item   self.\_send\_json(...) を実行する。
\item    を返す。
\item 条件 self.path.startswith('/metrics') を判定し、真なら内部処理を行う。
\item   self.\_send\_metrics(...) を実行する。
\item    を返す。
\item super().do\_GET(...) を実行する。
            \end{enumerate}
            \paragraph{代表コード断片}
            \begin{lstlisting}[language=Python]
    def do_GET(self):
        if self.path.startswith('/api/state'):
            self._send_json(self._node.get_state())
            return
        if self.path.startswith('/metrics'):
            self._send_metrics(self._node.get_prometheus_metrics())
            return
        super().do_GET()
\end{lstlisting}

\subsection{L33 関数 DashboardHandler.log\_message}
            \begin{itemize}[leftmargin=1.6em]
              \item 行範囲: L33--L35
              \item 所属: \path|DashboardHandler|
              \item 定義: \path|log_message(self, format, *args)|
              \item docstring: 明示されていない。
              \item このブロックが直接呼ぶ主な関数・メソッド: 特に明示されていない。
              \item 戻り値の要点: 戻り値が無いか、return 文が明示されていない。
              \item 主なローカル代入名: 特になし。
              \item 読み取る主なインスタンス属性: 特になし。
              \item 更新する主なインスタンス属性: 特になし。
              \item 明示的に送出する例外: 明示的raiseはない。
              \item 制御構造の規模: 条件分岐 0、ループ 0、try 0
            \end{itemize}
            \paragraph{この定義を読むためのPython構文}
            \begin{itemize}[leftmargin=1.6em]
            \item 第1引数selfは、このメソッドを呼んだインスタンス自身を受け取る慣習名である。
            \end{itemize}
            \paragraph{上から順の処理}
            \begin{enumerate}[leftmargin=1.8em]
            \item  を返す。
            \end{enumerate}
            \paragraph{代表コード断片}
            \begin{lstlisting}[language=Python]
    def log_message(self, format, *args):
        # quiet
        return
\end{lstlisting}

\subsection{L37 関数 DashboardHandler.\_send\_json}
            \begin{itemize}[leftmargin=1.6em]
              \item 行範囲: L37--L43
              \item 所属: \path|DashboardHandler|
              \item 定義: \path|_send_json(self, data)|
              \item docstring: 明示されていない。
              \item このブロックが直接呼ぶ主な関数・メソッド: dumps, encode, end\_headers, len, send\_header, send\_response, str, write
              \item 戻り値の要点: 戻り値が無いか、return 文が明示されていない。
              \item 主なローカル代入名: payload
              \item 読み取る主なインスタンス属性: self.end\_headers, self.send\_header, self.send\_response, self.wfile
              \item 更新する主なインスタンス属性: 特になし。
              \item 明示的に送出する例外: 明示的raiseはない。
              \item 制御構造の規模: 条件分岐 0、ループ 0、try 0
            \end{itemize}
            \paragraph{この定義を読むためのPython構文}
            \begin{itemize}[leftmargin=1.6em]
            \item 第1引数selfは、このメソッドを呼んだインスタンス自身を受け取る慣習名である。
            \end{itemize}
            \paragraph{上から順の処理}
            \begin{enumerate}[leftmargin=1.8em]
            \item payload に json.dumps(data).encode('utf-8') の結果を代入する。
\item self.send\_response(...) を実行する。
\item self.send\_header(...) を実行する。
\item self.send\_header(...) を実行する。
\item self.end\_headers(...) を実行する。
\item self.wfile.write(...) を実行する。
            \end{enumerate}
            \paragraph{代表コード断片}
            \begin{lstlisting}[language=Python]
    def _send_json(self, data):
        payload = json.dumps(data).encode('utf-8')
        self.send_response(200)
        self.send_header('Content-Type', 'application/json')
        self.send_header('Content-Length', str(len(payload)))
        self.end_headers()
        self.wfile.write(payload)
\end{lstlisting}

\subsection{L45 関数 DashboardHandler.\_send\_metrics}
            \begin{itemize}[leftmargin=1.6em]
              \item 行範囲: L45--L52
              \item 所属: \path|DashboardHandler|
              \item 定義: \path|_send_metrics(self, payload)|
              \item docstring: 明示されていない。
              \item このブロックが直接呼ぶ主な関数・メソッド: encode, end\_headers, len, send\_header, send\_response, str, write
              \item 戻り値の要点: 戻り値が無いか、return 文が明示されていない。
              \item 主なローカル代入名: body
              \item 読み取る主なインスタンス属性: self.end\_headers, self.send\_header, self.send\_response, self.wfile
              \item 更新する主なインスタンス属性: 特になし。
              \item 明示的に送出する例外: 明示的raiseはない。
              \item 制御構造の規模: 条件分岐 0、ループ 0、try 0
            \end{itemize}
            \paragraph{この定義を読むためのPython構文}
            \begin{itemize}[leftmargin=1.6em]
            \item 第1引数selfは、このメソッドを呼んだインスタンス自身を受け取る慣習名である。
            \end{itemize}
            \paragraph{上から順の処理}
            \begin{enumerate}[leftmargin=1.8em]
            \item body に payload.encode('utf-8') の結果を代入する。
\item self.send\_response(...) を実行する。
\item self.send\_header(...) を実行する。
\item self.send\_header(...) を実行する。
\item self.send\_header(...) を実行する。
\item self.end\_headers(...) を実行する。
\item self.wfile.write(...) を実行する。
            \end{enumerate}
            \paragraph{代表コード断片}
            \begin{lstlisting}[language=Python]
    def _send_metrics(self, payload):
        body = payload.encode('utf-8')
        self.send_response(200)
        self.send_header('Content-Type', 'text/plain; version=0.0.4; charset=utf-8')
        self.send_header('Cache-Control', 'no-store')
        self.send_header('Content-Length', str(len(body)))
        self.end_headers()
        self.wfile.write(body)
\end{lstlisting}

\subsection{L55 クラス DashboardNode}
            \begin{itemize}[leftmargin=1.6em]
              \item 行範囲: L55--L406
              \item 所属: module直下
              \item 定義: \path|DashboardNode(bases=Node)|
              \item docstring: 明示されていない。
              \item このブロックが直接呼ぶ主な関数・メソッド: 特に明示されていない。
              \item 戻り値の要点: 戻り値が無いか、return 文が明示されていない。
              \item 主なローカル代入名: 特になし。
              \item 読み取る主なインスタンス属性: 特になし。
              \item 更新する主なインスタンス属性: 特になし。
              \item 明示的に送出する例外: 明示的raiseはない。
              \item 制御構造の規模: 条件分岐 0、ループ 0、try 0
            \end{itemize}
            \paragraph{この定義を読むためのPython構文}
            \begin{itemize}[leftmargin=1.6em]
            \item class文は新しい型を定義する。丸括弧内は継承する基底クラスである。
\item 内包表記は、反復・条件・値生成を一つの式で記述してcollectionを作る。
\item with文はcontext managerに開始・終了処理を任せ、ファイルなどの資源を確実に解放する。
\item try文は例外が起き得る処理と、異常時または終了時の経路を分ける。
            \end{itemize}
            \paragraph{上から順の処理}
            \begin{enumerate}[leftmargin=1.8em]
            \item 関数 \_\_init\_\_ を定義する。
\item 関数 destroy\_node を定義する。
\item 関数 get\_state を定義する。
\item 関数 get\_prometheus\_metrics を定義する。
\item 関数 \_update を定義する。
\item 関数 \_on\_speed\_cmd を定義する。
\item 関数 \_on\_upper\_speed\_cmd を定義する。
\item 関数 \_on\_speed\_meas を定義する。
\item 関数 \_on\_throttle\_cmd を定義する。
\item 関数 \_on\_drive\_mode を定義する。
\item 関数 \_on\_soc を定義する。
\item 関数 \_set\_estimated\_soc を定義する。
            \end{enumerate}
            \paragraph{代表コード断片}
            \begin{lstlisting}[language=Python]
class DashboardNode(Node):
    def __init__(self):
        super().__init__('dashboard_node')
        self.declare_parameter('host', '0.0.0.0')
        self.declare_parameter('port', 8080)
        self.declare_parameter('static_dir', '')
        self.declare_parameter('dummy_csv', '')
        self.declare_parameter('dummy_rate_hz', 5.0)

        static_dir = str(self.get_parameter('static_dir').value).strip()
        if not static_dir:
            try:
                pkg_share = get_package_share_directory('mpc_solarcar')
                static_dir = os.path.join(pkg_share, 'dashboard')
            except Exception:
                static_dir = os.path.join(os.path.dirname(__file__), '..', 'dashboard')
        static_dir = os.path.abspath(static_dir)

        self._lock = threading.Lock()
        self._state = {
            'ts': time.time(),
            'speed_cmd_kmh': None,
            'upper_speed_cmd_kmh': None,
            'speed_meas_kmh': None,
            'model_speed_kmh': None,
            'throttle_cmd_pct': None,
            'drive_mode': None,
            'soc': None,
            'soc_meas': None,
            'soc_est': None,
            'Tb_C': None,
            's_km': None,
            'batt_current_a': None,
            'batt_voltage_v': None,
            'motor_w': None,
...
\end{lstlisting}

\subsection{L56 関数 DashboardNode.\_\_init\_\_}
            \begin{itemize}[leftmargin=1.6em]
              \item 行範囲: L56--L154
              \item 所属: \path|DashboardNode|
              \item 定義: \path|__init__(self)|
              \item docstring: 明示されていない。
              \item このブロックが直接呼ぶ主な関数・メソッド: Lock, Thread, ThreadingHTTPServer, \_\_init\_\_, \_init\_dummy, abspath, create\_subscription, declare\_parameter, dirname, error, float, get\_logger
              \item 戻り値の要点: 戻り値が無いか、return 文が明示されていない。
              \item 主なローカル代入名: dummy\_csv, handler, host, pkg\_share, port, static\_dir
              \item 読み取る主なインスタンス属性: self.\_init\_dummy, self.\_on\_drive\_mode, self.\_on\_env, self.\_on\_gps, self.\_on\_health, self.\_on\_ibatt, self.\_on\_lower\_plan, self.\_on\_metrics, self.\_on\_mpc\_state, self.\_on\_path, self.\_on\_s\_km, self.\_on\_soc, self.\_on\_speed\_cmd, self.\_on\_speed\_meas, self.\_on\_status, self.\_on\_sys\_diag, self.\_on\_sys\_state, self.\_on\_tb, self.\_on\_throttle\_cmd, self.\_on\_upper\_plan, self.\_on\_upper\_speed\_cmd, self.\_on\_vbatt, self.\_server, self.\_server\_thread
              \item 更新する主なインスタンス属性: self.\_lock, self.\_server, self.\_server\_thread, self.\_soc\_meas\_monotonic, self.\_state
              \item 明示的に送出する例外: 明示的raiseはない。
              \item 制御構造の規模: 条件分岐 2、ループ 0、try 2
            \end{itemize}
            \paragraph{この定義を読むためのPython構文}
            \begin{itemize}[leftmargin=1.6em]
            \item 第1引数selfは、このメソッドを呼んだインスタンス自身を受け取る慣習名である。
\item try文は例外が起き得る処理と、異常時または終了時の経路を分ける。
            \end{itemize}
            \paragraph{上から順の処理}
            \begin{enumerate}[leftmargin=1.8em]
            \item super().\_\_init\_\_(...) を実行する。
\item self.declare\_parameter(...) を実行する。
\item self.declare\_parameter(...) を実行する。
\item self.declare\_parameter(...) を実行する。
\item self.declare\_parameter(...) を実行する。
\item self.declare\_parameter(...) を実行する。
\item static\_dir に str(self.get\_parameter('static\_dir').value).strip() の結果を代入する。
\item 条件 not static\_dir を判定し、真なら内部処理を行う。
\item   例外処理を伴う try ブロックを実行する。
\item     pkg\_share に get\_package\_share\_directory('mpc\_solarcar') の結果を代入する。
\item     static\_dir に os.path.join(pkg\_share, 'dashboard') の結果を代入する。
\item     Exceptionを捕捉した場合:
\item     static\_dir に os.path.join(os.path.dirname(\_\_file\_\_), '..', 'dashboard') の結果を代入する。
\item static\_dir に os.path.abspath(static\_dir) の結果を代入する。
\item self.\_lock に threading.Lock() の結果を代入する。
\item self.\_state に \{'ts': time.time(), 'speed\_cmd\_kmh': None, 'upper\_speed\_cmd\_kmh': None, 'speed\_meas\_kmh': None, 'model\_speed\_kmh': None, 'throttle\_cmd\_pct': None, 'drive\_mode': None, 'soc': None, 'soc\_meas': None, 'soc\_est': None, 'Tb\_C': None, 's\_km': None, 'batt\_current\_a': None, 'batt\_voltage\_v': None, 'motor\_w': None, 'motor\_a': None, 'solar\_w': None, 'wheel\_w': None, 'pack\_w': None, 'G\_poa': None, 'Tcell\_C': None, 'Tamb\_C': None, 'headwind\_ms': None, 'slope\_pct': None, 'plan\_dt': None, 'lower\_dt': None, 'plan\_upper': None, 'plan\_lower': None, 'forecast\_k': None, 'sec\_to\_next': None, 'control\_stop\_hold': None, 'control\_stop\_remaining\_sec': None, 'control\_stop\_completed\_count': None, 'system\_state': None, 'system\_diag': None, 'mpc\_state': None, 'system\_health': None, 'gps\_lat': None, 'gps\_lon': None\} の結果を代入する。
\item self.\_soc\_meas\_monotonic に None の結果を代入する。
\item self.create\_subscription(...) を実行する。
\item self.create\_subscription(...) を実行する。
\item self.create\_subscription(...) を実行する。
\item self.create\_subscription(...) を実行する。
\item self.create\_subscription(...) を実行する。
\item self.create\_subscription(...) を実行する。
\item self.create\_subscription(...) を実行する。
\item self.create\_subscription(...) を実行する。
\item self.create\_subscription(...) を実行する。
\item self.create\_subscription(...) を実行する。
\item self.create\_subscription(...) を実行する。
\item self.create\_subscription(...) を実行する。
\item self.create\_subscription(...) を実行する。
\item self.create\_subscription(...) を実行する。
\item self.create\_subscription(...) を実行する。
\item self.create\_subscription(...) を実行する。
\item self.create\_subscription(...) を実行する。
\item self.create\_subscription(...) を実行する。
\item self.create\_subscription(...) を実行する。
\item self.create\_subscription(...) を実行する。
\item self.create\_subscription(...) を実行する。
\item dummy\_csv に str(self.get\_parameter('dummy\_csv').value).strip() の結果を代入する。
\item 条件 dummy\_csv を判定し、真なら内部処理を行う。
\item   self.\_init\_dummy(...) を実行する。
\item self.\_server に None の結果を代入する。
\item self.\_server\_thread に None の結果を代入する。
\item host に str(self.get\_parameter('host').value) の結果を代入する。
\item port に int(self.get\_parameter('port').value) の結果を代入する。
\item handler に partial(DashboardHandler, node=self, directory=static\_dir) の結果を代入する。
\item 例外処理を伴う try ブロックを実行する。
\item   self.\_server に ThreadingHTTPServer((host, port), handler) の結果を代入する。
\item   self.\_server\_thread に threading.Thread(target=self.\_server.serve\_forever, daemon=True) の結果を代入する。
\item   self.\_server\_thread.start(...) を実行する。
\item   self.get\_logger().info(...) を実行する。
\item   Exceptionを捕捉した場合:
\item   self.get\_logger().error(...) を実行する。
            \end{enumerate}
            \paragraph{代表コード断片}
            \begin{lstlisting}[language=Python]
    def __init__(self):
        super().__init__('dashboard_node')
        self.declare_parameter('host', '0.0.0.0')
        self.declare_parameter('port', 8080)
        self.declare_parameter('static_dir', '')
        self.declare_parameter('dummy_csv', '')
        self.declare_parameter('dummy_rate_hz', 5.0)

        static_dir = str(self.get_parameter('static_dir').value).strip()
        if not static_dir:
            try:
                pkg_share = get_package_share_directory('mpc_solarcar')
                static_dir = os.path.join(pkg_share, 'dashboard')
            except Exception:
                static_dir = os.path.join(os.path.dirname(__file__), '..', 'dashboard')
        static_dir = os.path.abspath(static_dir)

        self._lock = threading.Lock()
        self._state = {
            'ts': time.time(),
            'speed_cmd_kmh': None,
            'upper_speed_cmd_kmh': None,
            'speed_meas_kmh': None,
            'model_speed_kmh': None,
            'throttle_cmd_pct': None,
            'drive_mode': None,
            'soc': None,
            'soc_meas': None,
            'soc_est': None,
            'Tb_C': None,
            's_km': None,
            'batt_current_a': None,
            'batt_voltage_v': None,
            'motor_w': None,
            'motor_a': None,
...
\end{lstlisting}

\subsection{L156 関数 DashboardNode.destroy\_node}
            \begin{itemize}[leftmargin=1.6em]
              \item 行範囲: L156--L163
              \item 所属: \path|DashboardNode|
              \item 定義: \path|destroy_node(self)|
              \item docstring: 明示されていない。
              \item このブロックが直接呼ぶ主な関数・メソッド: destroy\_node, server\_close, shutdown, super
              \item 戻り値の要点: super().destroy\_node()
              \item 主なローカル代入名: 特になし。
              \item 読み取る主なインスタンス属性: self.\_server
              \item 更新する主なインスタンス属性: 特になし。
              \item 明示的に送出する例外: 明示的raiseはない。
              \item 制御構造の規模: 条件分岐 1、ループ 0、try 1
            \end{itemize}
            \paragraph{この定義を読むためのPython構文}
            \begin{itemize}[leftmargin=1.6em]
            \item 第1引数selfは、このメソッドを呼んだインスタンス自身を受け取る慣習名である。
\item try文は例外が起き得る処理と、異常時または終了時の経路を分ける。
            \end{itemize}
            \paragraph{上から順の処理}
            \begin{enumerate}[leftmargin=1.8em]
            \item 条件 self.\_server is not None を判定し、真なら内部処理を行う。
\item   例外処理を伴う try ブロックを実行する。
\item     self.\_server.shutdown(...) を実行する。
\item     self.\_server.server\_close(...) を実行する。
\item     Exceptionを捕捉した場合:
\item     Pass 文を実行する。
\item super().destroy\_node() を返す。
            \end{enumerate}
            \paragraph{代表コード断片}
            \begin{lstlisting}[language=Python]
    def destroy_node(self):
        if self._server is not None:
            try:
                self._server.shutdown()
                self._server.server_close()
            except Exception:
                pass
        return super().destroy_node()
\end{lstlisting}

\subsection{L165 関数 DashboardNode.get\_state}
            \begin{itemize}[leftmargin=1.6em]
              \item 行範囲: L165--L167
              \item 所属: \path|DashboardNode|
              \item 定義: \path|get_state(self)|
              \item docstring: 明示されていない。
              \item このブロックが直接呼ぶ主な関数・メソッド: dict
              \item 戻り値の要点: dict(self.\_state)
              \item 主なローカル代入名: 特になし。
              \item 読み取る主なインスタンス属性: self.\_lock, self.\_state
              \item 更新する主なインスタンス属性: 特になし。
              \item 明示的に送出する例外: 明示的raiseはない。
              \item 制御構造の規模: 条件分岐 0、ループ 0、try 0
            \end{itemize}
            \paragraph{この定義を読むためのPython構文}
            \begin{itemize}[leftmargin=1.6em]
            \item 第1引数selfは、このメソッドを呼んだインスタンス自身を受け取る慣習名である。
\item with文はcontext managerに開始・終了処理を任せ、ファイルなどの資源を確実に解放する。
            \end{itemize}
            \paragraph{上から順の処理}
            \begin{enumerate}[leftmargin=1.8em]
            \item with 文で self.\_lock を管理しながら処理する。
\item   dict(self.\_state) を返す。
            \end{enumerate}
            \paragraph{代表コード断片}
            \begin{lstlisting}[language=Python]
    def get_state(self):
        with self._lock:
            return dict(self._state)
\end{lstlisting}

\subsection{L169 関数 DashboardNode.get\_prometheus\_metrics}
            \begin{itemize}[leftmargin=1.6em]
              \item 行範囲: L169--L234
              \item 所属: \path|DashboardNode|
              \item 定義: \path|get_prometheus_metrics(self)|
              \item docstring: 明示されていない。
              \item このブロックが直接呼ぶ主な関数・メソッド: dict, extend, float, get, isfinite, items, join, max, replace, str, time
              \item 戻り値の要点: '\textbackslash{}n'.join(lines) + '\textbackslash{}n'
              \item 主なローカル代入名: age\_sec, help\_text, key, label, lines, metric\_keys, metric\_name, number, state, state\_key, value
              \item 読み取る主なインスタンス属性: self.\_lock, self.\_state
              \item 更新する主なインスタンス属性: 特になし。
              \item 明示的に送出する例外: 明示的raiseはない。
              \item 制御構造の規模: 条件分岐 1、ループ 2、try 1
            \end{itemize}
            \paragraph{この定義を読むためのPython構文}
            \begin{itemize}[leftmargin=1.6em]
            \item 第1引数selfは、このメソッドを呼んだインスタンス自身を受け取る慣習名である。
\item with文はcontext managerに開始・終了処理を任せ、ファイルなどの資源を確実に解放する。
\item try文は例外が起き得る処理と、異常時または終了時の経路を分ける。
            \end{itemize}
            \paragraph{上から順の処理}
            \begin{enumerate}[leftmargin=1.8em]
            \item metric\_keys に \{'speed\_cmd\_kmh': ('solarcar\_speed\_command\_kmh', 'Lower speed command'), 'upper\_speed\_cmd\_kmh': ('solarcar\_speed\_upper\_command\_kmh', 'Upper planner speed command'), 'speed\_meas\_kmh': ('solarcar\_speed\_measured\_kmh', 'Measured vehicle speed'), 'model\_speed\_kmh': ('solarcar\_speed\_model\_kmh', 'Speed used for planner model metrics'), 'throttle\_cmd\_pct': ('solarcar\_throttle\_command\_percent', 'Lower controller throttle command'), 'soc': ('solarcar\_battery\_soc\_ratio', 'Selected battery state of charge'), 'soc\_meas': ('solarcar\_battery\_soc\_measured\_ratio', 'Measured battery state of charge'), 'soc\_est': ('solarcar\_battery\_soc\_estimated\_ratio', 'Estimated battery state of charge'), 'Tb\_C': ('solarcar\_battery\_temperature\_celsius', 'Battery temperature'), 's\_km': ('solarcar\_distance\_km', 'Race distance'), 'batt\_current\_a': ('solarcar\_battery\_current\_ampere', 'Battery current'), 'batt\_voltage\_v': ('solarcar\_battery\_voltage\_volt', 'Battery terminal voltage'), 'motor\_w': ('solarcar\_motor\_power\_watt', 'Motor electrical power'), 'motor\_a': ('solarcar\_motor\_current\_ampere', 'Motor current'), 'solar\_w': ('solarcar\_pv\_power\_watt', 'PV power'), 'wheel\_w': ('solarcar\_wheel\_power\_watt', 'Wheel mechanical power'), 'pack\_w': ('solarcar\_pack\_power\_watt', 'Battery pack power'), 'G\_poa': ('solarcar\_irradiance\_poa\_watt\_per\_m2', 'Plane-of-array irradiance'), 'Tcell\_C': ('solarcar\_cell\_temperature\_celsius', 'PV cell temperature'), 'Tamb\_C': ('solarcar\_ambient\_temperature\_celsius', 'Ambient temperature'), 'headwind\_ms': ('solarcar\_headwind\_meter\_per\_second', 'Corrected headwind'), 'slope\_pct': ('solarcar\_route\_slope\_percent', 'Route slope'), 'forecast\_k': ('solarcar\_forecast\_index', 'Forecast row index'), 'sec\_to\_next': ('solarcar\_seconds\_to\_next\_update', 'Seconds to next planner update'), 'control\_stop\_hold': ('solarcar\_control\_stop\_hold', 'Control-stop approach or dwell is active'), 'control\_stop\_remaining\_sec': ('solarcar\_control\_stop\_remaining\_seconds', 'Remaining mandatory control-stop dwell'), 'control\_stop\_completed\_count': ('solarcar\_control\_stop\_completed\_total', 'Completed mandatory control stops'), 'system\_health': ('solarcar\_system\_health\_ratio', 'System health ratio'), 'gps\_lat': ('solarcar\_gps\_latitude\_degree', 'GNSS latitude'), 'gps\_lon': ('solarcar\_gps\_longitude\_degree', 'GNSS longitude'), 'path\_points': ('solarcar\_path\_points', 'Published trajectory point count')\} の結果を代入する。
\item with 文で self.\_lock を管理しながら処理する。
\item   state に dict(self.\_state) の結果を代入する。
\item lines に [] の結果を代入する。
\item metric\_keys.items() を順に走査し、各要素を (state\_key, (metric\_name, help\_text)) に入れて処理する。
\item   value に state.get(state\_key) の結果を代入する。
\item   例外処理を伴う try ブロックを実行する。
\item     number に float(value) の結果を代入する。
\item     (TypeError, ValueError)を捕捉した場合:
\item     Continue 文を実行する。
\item   条件 not math.isfinite(number) を判定し、真なら内部処理を行う。
\item     Continue 文を実行する。
\item   lines.extend(...) を実行する。
\item age\_sec に max(0.0, time.time() - float(state.get('ts', time.time()))) の結果を代入する。
\item lines.extend(...) を実行する。
\item (('system\_state', 'solarcar\_system\_state\_info'), ('mpc\_state', 'solarcar\_mpc\_state\_info'), ('drive\_mode', 'solarcar\_drive\_mode\_info')) を順に走査し、各要素を (key, metric\_name) に入れて処理する。
\item   label に str(state.get(key) or 'unknown').replace('\textbackslash{}\textbackslash{}', '\textbackslash{}\textbackslash{}\textbackslash{}\textbackslash{}').replace('"', '\textbackslash{}\textbackslash{}"').replace('\textbackslash{}n', ' ') の結果を代入する。
\item   lines.extend(...) を実行する。
\item '\textbackslash{}n'.join(lines) + '\textbackslash{}n' を返す。
            \end{enumerate}
            \paragraph{代表コード断片}
            \begin{lstlisting}[language=Python]
    def get_prometheus_metrics(self):
        metric_keys = {
            'speed_cmd_kmh': ('solarcar_speed_command_kmh', 'Lower speed command'),
            'upper_speed_cmd_kmh': ('solarcar_speed_upper_command_kmh', 'Upper planner speed command'),
            'speed_meas_kmh': ('solarcar_speed_measured_kmh', 'Measured vehicle speed'),
            'model_speed_kmh': ('solarcar_speed_model_kmh', 'Speed used for planner model metrics'),
            'throttle_cmd_pct': ('solarcar_throttle_command_percent', 'Lower controller throttle command'),
            'soc': ('solarcar_battery_soc_ratio', 'Selected battery state of charge'),
            'soc_meas': ('solarcar_battery_soc_measured_ratio', 'Measured battery state of charge'),
            'soc_est': ('solarcar_battery_soc_estimated_ratio', 'Estimated battery state of charge'),
            'Tb_C': ('solarcar_battery_temperature_celsius', 'Battery temperature'),
            's_km': ('solarcar_distance_km', 'Race distance'),
            'batt_current_a': ('solarcar_battery_current_ampere', 'Battery current'),
            'batt_voltage_v': ('solarcar_battery_voltage_volt', 'Battery terminal voltage'),
            'motor_w': ('solarcar_motor_power_watt', 'Motor electrical power'),
            'motor_a': ('solarcar_motor_current_ampere', 'Motor current'),
            'solar_w': ('solarcar_pv_power_watt', 'PV power'),
            'wheel_w': ('solarcar_wheel_power_watt', 'Wheel mechanical power'),
            'pack_w': ('solarcar_pack_power_watt', 'Battery pack power'),
            'G_poa': ('solarcar_irradiance_poa_watt_per_m2', 'Plane-of-array irradiance'),
            'Tcell_C': ('solarcar_cell_temperature_celsius', 'PV cell temperature'),
            'Tamb_C': ('solarcar_ambient_temperature_celsius', 'Ambient temperature'),
            'headwind_ms': ('solarcar_headwind_meter_per_second', 'Corrected headwind'),
            'slope_pct': ('solarcar_route_slope_percent', 'Route slope'),
            'forecast_k': ('solarcar_forecast_index', 'Forecast row index'),
            'sec_to_next': ('solarcar_seconds_to_next_update', 'Seconds to next planner update'),
            'control_stop_hold': ('solarcar_control_stop_hold', 'Control-stop approach or dwell is active'),
            'control_stop_remaining_sec': ('solarcar_control_stop_remaining_seconds', 'Remaining mandatory control-stop dwell'),
            'control_stop_completed_count': ('solarcar_control_stop_completed_total', 'Completed mandatory control stops'),
            'system_health': ('solarcar_system_health_ratio', 'System health ratio'),
            'gps_lat': ('solarcar_gps_latitude_degree', 'GNSS latitude'),
            'gps_lon': ('solarcar_gps_longitude_degree', 'GNSS longitude'),
            'path_points': ('solarcar_path_points', 'Published trajectory point count'),
        }
        with self._lock:
...
\end{lstlisting}

\subsection{L236 関数 DashboardNode.\_update}
            \begin{itemize}[leftmargin=1.6em]
              \item 行範囲: L236--L239
              \item 所属: \path|DashboardNode|
              \item 定義: \path|_update(self, key, value)|
              \item docstring: 明示されていない。
              \item このブロックが直接呼ぶ主な関数・メソッド: time
              \item 戻り値の要点: 戻り値が無いか、return 文が明示されていない。
              \item 主なローカル代入名: 特になし。
              \item 読み取る主なインスタンス属性: self.\_lock, self.\_state
              \item 更新する主なインスタンス属性: 特になし。
              \item 明示的に送出する例外: 明示的raiseはない。
              \item 制御構造の規模: 条件分岐 0、ループ 0、try 0
            \end{itemize}
            \paragraph{この定義を読むためのPython構文}
            \begin{itemize}[leftmargin=1.6em]
            \item 第1引数selfは、このメソッドを呼んだインスタンス自身を受け取る慣習名である。
\item with文はcontext managerに開始・終了処理を任せ、ファイルなどの資源を確実に解放する。
            \end{itemize}
            \paragraph{上から順の処理}
            \begin{enumerate}[leftmargin=1.8em]
            \item with 文で self.\_lock を管理しながら処理する。
\item   self.\_state[key] に value の結果を代入する。
\item   self.\_state['ts'] に time.time() の結果を代入する。
            \end{enumerate}
            \paragraph{代表コード断片}
            \begin{lstlisting}[language=Python]
    def _update(self, key, value):
        with self._lock:
            self._state[key] = value
            self._state['ts'] = time.time()
\end{lstlisting}

\subsection{L241 関数 DashboardNode.\_on\_speed\_cmd}
            \begin{itemize}[leftmargin=1.6em]
              \item 行範囲: L241--L242
              \item 所属: \path|DashboardNode|
              \item 定義: \path|_on_speed_cmd(self, msg: Float32)|
              \item docstring: 明示されていない。
              \item このブロックが直接呼ぶ主な関数・メソッド: \_update, float
              \item 戻り値の要点: 戻り値が無いか、return 文が明示されていない。
              \item 主なローカル代入名: 特になし。
              \item 読み取る主なインスタンス属性: self.\_update
              \item 更新する主なインスタンス属性: 特になし。
              \item 明示的に送出する例外: 明示的raiseはない。
              \item 制御構造の規模: 条件分岐 0、ループ 0、try 0
            \end{itemize}
            \paragraph{この定義を読むためのPython構文}
            \begin{itemize}[leftmargin=1.6em]
            \item 第1引数selfは、このメソッドを呼んだインスタンス自身を受け取る慣習名である。
            \end{itemize}
            \paragraph{上から順の処理}
            \begin{enumerate}[leftmargin=1.8em]
            \item self.\_update(...) を実行する。
            \end{enumerate}
            \paragraph{代表コード断片}
            \begin{lstlisting}[language=Python]
    def _on_speed_cmd(self, msg: Float32):
        self._update('speed_cmd_kmh', float(msg.data))
\end{lstlisting}

\subsection{L244 関数 DashboardNode.\_on\_upper\_speed\_cmd}
            \begin{itemize}[leftmargin=1.6em]
              \item 行範囲: L244--L245
              \item 所属: \path|DashboardNode|
              \item 定義: \path|_on_upper_speed_cmd(self, msg: Float32)|
              \item docstring: 明示されていない。
              \item このブロックが直接呼ぶ主な関数・メソッド: \_update, float
              \item 戻り値の要点: 戻り値が無いか、return 文が明示されていない。
              \item 主なローカル代入名: 特になし。
              \item 読み取る主なインスタンス属性: self.\_update
              \item 更新する主なインスタンス属性: 特になし。
              \item 明示的に送出する例外: 明示的raiseはない。
              \item 制御構造の規模: 条件分岐 0、ループ 0、try 0
            \end{itemize}
            \paragraph{この定義を読むためのPython構文}
            \begin{itemize}[leftmargin=1.6em]
            \item 第1引数selfは、このメソッドを呼んだインスタンス自身を受け取る慣習名である。
            \end{itemize}
            \paragraph{上から順の処理}
            \begin{enumerate}[leftmargin=1.8em]
            \item self.\_update(...) を実行する。
            \end{enumerate}
            \paragraph{代表コード断片}
            \begin{lstlisting}[language=Python]
    def _on_upper_speed_cmd(self, msg: Float32):
        self._update('upper_speed_cmd_kmh', float(msg.data))
\end{lstlisting}

\subsection{L247 関数 DashboardNode.\_on\_speed\_meas}
            \begin{itemize}[leftmargin=1.6em]
              \item 行範囲: L247--L248
              \item 所属: \path|DashboardNode|
              \item 定義: \path|_on_speed_meas(self, msg: Float32)|
              \item docstring: 明示されていない。
              \item このブロックが直接呼ぶ主な関数・メソッド: \_update, float
              \item 戻り値の要点: 戻り値が無いか、return 文が明示されていない。
              \item 主なローカル代入名: 特になし。
              \item 読み取る主なインスタンス属性: self.\_update
              \item 更新する主なインスタンス属性: 特になし。
              \item 明示的に送出する例外: 明示的raiseはない。
              \item 制御構造の規模: 条件分岐 0、ループ 0、try 0
            \end{itemize}
            \paragraph{この定義を読むためのPython構文}
            \begin{itemize}[leftmargin=1.6em]
            \item 第1引数selfは、このメソッドを呼んだインスタンス自身を受け取る慣習名である。
            \end{itemize}
            \paragraph{上から順の処理}
            \begin{enumerate}[leftmargin=1.8em]
            \item self.\_update(...) を実行する。
            \end{enumerate}
            \paragraph{代表コード断片}
            \begin{lstlisting}[language=Python]
    def _on_speed_meas(self, msg: Float32):
        self._update('speed_meas_kmh', float(msg.data))
\end{lstlisting}

\subsection{L250 関数 DashboardNode.\_on\_throttle\_cmd}
            \begin{itemize}[leftmargin=1.6em]
              \item 行範囲: L250--L251
              \item 所属: \path|DashboardNode|
              \item 定義: \path|_on_throttle_cmd(self, msg: Float32)|
              \item docstring: 明示されていない。
              \item このブロックが直接呼ぶ主な関数・メソッド: \_update, float
              \item 戻り値の要点: 戻り値が無いか、return 文が明示されていない。
              \item 主なローカル代入名: 特になし。
              \item 読み取る主なインスタンス属性: self.\_update
              \item 更新する主なインスタンス属性: 特になし。
              \item 明示的に送出する例外: 明示的raiseはない。
              \item 制御構造の規模: 条件分岐 0、ループ 0、try 0
            \end{itemize}
            \paragraph{この定義を読むためのPython構文}
            \begin{itemize}[leftmargin=1.6em]
            \item 第1引数selfは、このメソッドを呼んだインスタンス自身を受け取る慣習名である。
            \end{itemize}
            \paragraph{上から順の処理}
            \begin{enumerate}[leftmargin=1.8em]
            \item self.\_update(...) を実行する。
            \end{enumerate}
            \paragraph{代表コード断片}
            \begin{lstlisting}[language=Python]
    def _on_throttle_cmd(self, msg: Float32):
        self._update('throttle_cmd_pct', float(msg.data))
\end{lstlisting}

\subsection{L253 関数 DashboardNode.\_on\_drive\_mode}
            \begin{itemize}[leftmargin=1.6em]
              \item 行範囲: L253--L254
              \item 所属: \path|DashboardNode|
              \item 定義: \path|_on_drive_mode(self, msg: String)|
              \item docstring: 明示されていない。
              \item このブロックが直接呼ぶ主な関数・メソッド: \_update, str
              \item 戻り値の要点: 戻り値が無いか、return 文が明示されていない。
              \item 主なローカル代入名: 特になし。
              \item 読み取る主なインスタンス属性: self.\_update
              \item 更新する主なインスタンス属性: 特になし。
              \item 明示的に送出する例外: 明示的raiseはない。
              \item 制御構造の規模: 条件分岐 0、ループ 0、try 0
            \end{itemize}
            \paragraph{この定義を読むためのPython構文}
            \begin{itemize}[leftmargin=1.6em]
            \item 第1引数selfは、このメソッドを呼んだインスタンス自身を受け取る慣習名である。
            \end{itemize}
            \paragraph{上から順の処理}
            \begin{enumerate}[leftmargin=1.8em]
            \item self.\_update(...) を実行する。
            \end{enumerate}
            \paragraph{代表コード断片}
            \begin{lstlisting}[language=Python]
    def _on_drive_mode(self, msg: String):
        self._update('drive_mode', str(msg.data))
\end{lstlisting}

\subsection{L256 関数 DashboardNode.\_on\_soc}
            \begin{itemize}[leftmargin=1.6em]
              \item 行範囲: L256--L262
              \item 所属: \path|DashboardNode|
              \item 定義: \path|_on_soc(self, msg: Float32)|
              \item docstring: 明示されていない。
              \item このブロックが直接呼ぶ主な関数・メソッド: float, monotonic, time
              \item 戻り値の要点: 戻り値が無いか、return 文が明示されていない。
              \item 主なローカル代入名: value
              \item 読み取る主なインスタンス属性: self.\_lock, self.\_state
              \item 更新する主なインスタンス属性: self.\_soc\_meas\_monotonic
              \item 明示的に送出する例外: 明示的raiseはない。
              \item 制御構造の規模: 条件分岐 0、ループ 0、try 0
            \end{itemize}
            \paragraph{この定義を読むためのPython構文}
            \begin{itemize}[leftmargin=1.6em]
            \item 第1引数selfは、このメソッドを呼んだインスタンス自身を受け取る慣習名である。
\item with文はcontext managerに開始・終了処理を任せ、ファイルなどの資源を確実に解放する。
            \end{itemize}
            \paragraph{上から順の処理}
            \begin{enumerate}[leftmargin=1.8em]
            \item value に float(msg.data) の結果を代入する。
\item with 文で self.\_lock を管理しながら処理する。
\item   self.\_state['soc'] に value の結果を代入する。
\item   self.\_state['soc\_meas'] に value の結果を代入する。
\item   self.\_state['ts'] に time.time() の結果を代入する。
\item   self.\_soc\_meas\_monotonic に time.monotonic() の結果を代入する。
            \end{enumerate}
            \paragraph{代表コード断片}
            \begin{lstlisting}[language=Python]
    def _on_soc(self, msg: Float32):
        value = float(msg.data)
        with self._lock:
            self._state['soc'] = value
            self._state['soc_meas'] = value
            self._state['ts'] = time.time()
            self._soc_meas_monotonic = time.monotonic()
\end{lstlisting}

\subsection{L264 関数 DashboardNode.\_set\_estimated\_soc}
            \begin{itemize}[leftmargin=1.6em]
              \item 行範囲: L264--L271
              \item 所属: \path|DashboardNode|
              \item 定義: \path|_set_estimated_soc(self, value)|
              \item docstring: 明示されていない。
              \item このブロックが直接呼ぶ主な関数・メソッド: float, monotonic
              \item 戻り値の要点: 戻り値が無いか、return 文が明示されていない。
              \item 主なローカル代入名: measured\_is\_fresh
              \item 読み取る主なインスタンス属性: self.\_soc\_meas\_monotonic, self.\_state
              \item 更新する主なインスタンス属性: 特になし。
              \item 明示的に送出する例外: 明示的raiseはない。
              \item 制御構造の規模: 条件分岐 1、ループ 0、try 0
            \end{itemize}
            \paragraph{この定義を読むためのPython構文}
            \begin{itemize}[leftmargin=1.6em]
            \item 第1引数selfは、このメソッドを呼んだインスタンス自身を受け取る慣習名である。
            \end{itemize}
            \paragraph{上から順の処理}
            \begin{enumerate}[leftmargin=1.8em]
            \item self.\_state['soc\_est'] に float(value) の結果を代入する。
\item measured\_is\_fresh に self.\_soc\_meas\_monotonic is not None and time.monotonic() - self.\_soc\_meas\_monotonic <= 5.0 の結果を代入する。
\item 条件 not measured\_is\_fresh を判定し、真なら内部処理を行う。
\item   self.\_state['soc'] に float(value) の結果を代入する。
            \end{enumerate}
            \paragraph{代表コード断片}
            \begin{lstlisting}[language=Python]
    def _set_estimated_soc(self, value):
        self._state['soc_est'] = float(value)
        measured_is_fresh = (
            self._soc_meas_monotonic is not None
            and (time.monotonic() - self._soc_meas_monotonic) <= 5.0
        )
        if not measured_is_fresh:
            self._state['soc'] = float(value)
\end{lstlisting}

\subsection{L273 関数 DashboardNode.\_on\_tb}
            \begin{itemize}[leftmargin=1.6em]
              \item 行範囲: L273--L274
              \item 所属: \path|DashboardNode|
              \item 定義: \path|_on_tb(self, msg: Float32)|
              \item docstring: 明示されていない。
              \item このブロックが直接呼ぶ主な関数・メソッド: \_update, float
              \item 戻り値の要点: 戻り値が無いか、return 文が明示されていない。
              \item 主なローカル代入名: 特になし。
              \item 読み取る主なインスタンス属性: self.\_update
              \item 更新する主なインスタンス属性: 特になし。
              \item 明示的に送出する例外: 明示的raiseはない。
              \item 制御構造の規模: 条件分岐 0、ループ 0、try 0
            \end{itemize}
            \paragraph{この定義を読むためのPython構文}
            \begin{itemize}[leftmargin=1.6em]
            \item 第1引数selfは、このメソッドを呼んだインスタンス自身を受け取る慣習名である。
            \end{itemize}
            \paragraph{上から順の処理}
            \begin{enumerate}[leftmargin=1.8em]
            \item self.\_update(...) を実行する。
            \end{enumerate}
            \paragraph{代表コード断片}
            \begin{lstlisting}[language=Python]
    def _on_tb(self, msg: Float32):
        self._update('Tb_C', float(msg.data))
\end{lstlisting}

\subsection{L276 関数 DashboardNode.\_on\_ibatt}
            \begin{itemize}[leftmargin=1.6em]
              \item 行範囲: L276--L277
              \item 所属: \path|DashboardNode|
              \item 定義: \path|_on_ibatt(self, msg: Float32)|
              \item docstring: 明示されていない。
              \item このブロックが直接呼ぶ主な関数・メソッド: \_update, float
              \item 戻り値の要点: 戻り値が無いか、return 文が明示されていない。
              \item 主なローカル代入名: 特になし。
              \item 読み取る主なインスタンス属性: self.\_update
              \item 更新する主なインスタンス属性: 特になし。
              \item 明示的に送出する例外: 明示的raiseはない。
              \item 制御構造の規模: 条件分岐 0、ループ 0、try 0
            \end{itemize}
            \paragraph{この定義を読むためのPython構文}
            \begin{itemize}[leftmargin=1.6em]
            \item 第1引数selfは、このメソッドを呼んだインスタンス自身を受け取る慣習名である。
            \end{itemize}
            \paragraph{上から順の処理}
            \begin{enumerate}[leftmargin=1.8em]
            \item self.\_update(...) を実行する。
            \end{enumerate}
            \paragraph{代表コード断片}
            \begin{lstlisting}[language=Python]
    def _on_ibatt(self, msg: Float32):
        self._update('batt_current_a', float(msg.data))
\end{lstlisting}

\subsection{L279 関数 DashboardNode.\_on\_vbatt}
            \begin{itemize}[leftmargin=1.6em]
              \item 行範囲: L279--L280
              \item 所属: \path|DashboardNode|
              \item 定義: \path|_on_vbatt(self, msg: Float32)|
              \item docstring: 明示されていない。
              \item このブロックが直接呼ぶ主な関数・メソッド: \_update, float
              \item 戻り値の要点: 戻り値が無いか、return 文が明示されていない。
              \item 主なローカル代入名: 特になし。
              \item 読み取る主なインスタンス属性: self.\_update
              \item 更新する主なインスタンス属性: 特になし。
              \item 明示的に送出する例外: 明示的raiseはない。
              \item 制御構造の規模: 条件分岐 0、ループ 0、try 0
            \end{itemize}
            \paragraph{この定義を読むためのPython構文}
            \begin{itemize}[leftmargin=1.6em]
            \item 第1引数selfは、このメソッドを呼んだインスタンス自身を受け取る慣習名である。
            \end{itemize}
            \paragraph{上から順の処理}
            \begin{enumerate}[leftmargin=1.8em]
            \item self.\_update(...) を実行する。
            \end{enumerate}
            \paragraph{代表コード断片}
            \begin{lstlisting}[language=Python]
    def _on_vbatt(self, msg: Float32):
        self._update('batt_voltage_v', float(msg.data))
\end{lstlisting}

\subsection{L282 関数 DashboardNode.\_on\_s\_km}
            \begin{itemize}[leftmargin=1.6em]
              \item 行範囲: L282--L283
              \item 所属: \path|DashboardNode|
              \item 定義: \path|_on_s_km(self, msg: Float32)|
              \item docstring: 明示されていない。
              \item このブロックが直接呼ぶ主な関数・メソッド: \_update, float
              \item 戻り値の要点: 戻り値が無いか、return 文が明示されていない。
              \item 主なローカル代入名: 特になし。
              \item 読み取る主なインスタンス属性: self.\_update
              \item 更新する主なインスタンス属性: 特になし。
              \item 明示的に送出する例外: 明示的raiseはない。
              \item 制御構造の規模: 条件分岐 0、ループ 0、try 0
            \end{itemize}
            \paragraph{この定義を読むためのPython構文}
            \begin{itemize}[leftmargin=1.6em]
            \item 第1引数selfは、このメソッドを呼んだインスタンス自身を受け取る慣習名である。
            \end{itemize}
            \paragraph{上から順の処理}
            \begin{enumerate}[leftmargin=1.8em]
            \item self.\_update(...) を実行する。
            \end{enumerate}
            \paragraph{代表コード断片}
            \begin{lstlisting}[language=Python]
    def _on_s_km(self, msg: Float32):
        self._update('s_km', float(msg.data))
\end{lstlisting}

\subsection{L285 関数 DashboardNode.\_on\_status}
            \begin{itemize}[leftmargin=1.6em]
              \item 行範囲: L285--L299
              \item 所属: \path|DashboardNode|
              \item 定義: \path|_on_status(self, msg: Float32MultiArray)|
              \item docstring: 明示されていない。
              \item このブロックが直接呼ぶ主な関数・メソッド: \_set\_estimated\_soc, float, len, list, time
              \item 戻り値の要点: 戻り値が無いか、return 文が明示されていない。
              \item 主なローカル代入名: data
              \item 読み取る主なインスタンス属性: self.\_lock, self.\_set\_estimated\_soc, self.\_state
              \item 更新する主なインスタンス属性: 特になし。
              \item 明示的に送出する例外: 明示的raiseはない。
              \item 制御構造の規模: 条件分岐 2、ループ 0、try 0
            \end{itemize}
            \paragraph{この定義を読むためのPython構文}
            \begin{itemize}[leftmargin=1.6em]
            \item 第1引数selfは、このメソッドを呼んだインスタンス自身を受け取る慣習名である。
\item with文はcontext managerに開始・終了処理を任せ、ファイルなどの資源を確実に解放する。
            \end{itemize}
            \paragraph{上から順の処理}
            \begin{enumerate}[leftmargin=1.8em]
            \item data に list(msg.data) の結果を代入する。
\item with 文で self.\_lock を管理しながら処理する。
\item   条件 len(data) >= 5 を判定し、真なら内部処理を行う。
\item     self.\_set\_estimated\_soc(...) を実行する。
\item     self.\_state['Tb\_C'] に float(data[1]) の結果を代入する。
\item     self.\_state['s\_km'] に float(data[2]) の結果を代入する。
\item     self.\_state['forecast\_k'] に float(data[3]) の結果を代入する。
\item     self.\_state['sec\_to\_next'] に float(data[4]) の結果を代入する。
\item   条件 len(data) >= 8 を判定し、真なら内部処理を行う。
\item     self.\_state['control\_stop\_hold'] に float(data[5]) の結果を代入する。
\item     self.\_state['control\_stop\_remaining\_sec'] に float(data[6]) の結果を代入する。
\item     self.\_state['control\_stop\_completed\_count'] に float(data[7]) の結果を代入する。
\item   self.\_state['status\_raw'] に data の結果を代入する。
\item   self.\_state['ts'] に time.time() の結果を代入する。
            \end{enumerate}
            \paragraph{代表コード断片}
            \begin{lstlisting}[language=Python]
    def _on_status(self, msg: Float32MultiArray):
        data = list(msg.data)
        with self._lock:
            if len(data) >= 5:
                self._set_estimated_soc(data[0])
                self._state['Tb_C'] = float(data[1])
                self._state['s_km'] = float(data[2])
                self._state['forecast_k'] = float(data[3])
                self._state['sec_to_next'] = float(data[4])
            if len(data) >= 8:
                self._state['control_stop_hold'] = float(data[5])
                self._state['control_stop_remaining_sec'] = float(data[6])
                self._state['control_stop_completed_count'] = float(data[7])
            self._state['status_raw'] = data
            self._state['ts'] = time.time()
\end{lstlisting}

\subsection{L301 関数 DashboardNode.\_on\_env}
            \begin{itemize}[leftmargin=1.6em]
              \item 行範囲: L301--L311
              \item 所属: \path|DashboardNode|
              \item 定義: \path|_on_env(self, msg: Float32MultiArray)|
              \item docstring: 明示されていない。
              \item このブロックが直接呼ぶ主な関数・メソッド: float, len, list, time
              \item 戻り値の要点: 戻り値が無いか、return 文が明示されていない。
              \item 主なローカル代入名: data
              \item 読み取る主なインスタンス属性: self.\_lock, self.\_state
              \item 更新する主なインスタンス属性: 特になし。
              \item 明示的に送出する例外: 明示的raiseはない。
              \item 制御構造の規模: 条件分岐 1、ループ 0、try 0
            \end{itemize}
            \paragraph{この定義を読むためのPython構文}
            \begin{itemize}[leftmargin=1.6em]
            \item 第1引数selfは、このメソッドを呼んだインスタンス自身を受け取る慣習名である。
\item with文はcontext managerに開始・終了処理を任せ、ファイルなどの資源を確実に解放する。
            \end{itemize}
            \paragraph{上から順の処理}
            \begin{enumerate}[leftmargin=1.8em]
            \item data に list(msg.data) の結果を代入する。
\item 条件 len(data) < 5 を判定し、真なら内部処理を行う。
\item    を返す。
\item with 文で self.\_lock を管理しながら処理する。
\item   self.\_state['G\_poa'] に float(data[0]) の結果を代入する。
\item   self.\_state['Tcell\_C'] に float(data[1]) の結果を代入する。
\item   self.\_state['Tamb\_C'] に float(data[2]) の結果を代入する。
\item   self.\_state['slope\_pct'] に float(data[3]) の結果を代入する。
\item   self.\_state['headwind\_ms'] に float(data[4]) の結果を代入する。
\item   self.\_state['ts'] に time.time() の結果を代入する。
            \end{enumerate}
            \paragraph{代表コード断片}
            \begin{lstlisting}[language=Python]
    def _on_env(self, msg: Float32MultiArray):
        data = list(msg.data)
        if len(data) < 5:
            return
        with self._lock:
            self._state['G_poa'] = float(data[0])
            self._state['Tcell_C'] = float(data[1])
            self._state['Tamb_C'] = float(data[2])
            self._state['slope_pct'] = float(data[3])
            self._state['headwind_ms'] = float(data[4])
            self._state['ts'] = time.time()
\end{lstlisting}

\subsection{L313 関数 DashboardNode.\_on\_metrics}
            \begin{itemize}[leftmargin=1.6em]
              \item 行範囲: L313--L331
              \item 所属: \path|DashboardNode|
              \item 定義: \path|_on_metrics(self, msg: Float32MultiArray)|
              \item docstring: 明示されていない。
              \item このブロックが直接呼ぶ主な関数・メソッド: \_set\_estimated\_soc, float, len, list, time
              \item 戻り値の要点: 戻り値が無いか、return 文が明示されていない。
              \item 主なローカル代入名: data
              \item 読み取る主なインスタンス属性: self.\_lock, self.\_set\_estimated\_soc, self.\_state
              \item 更新する主なインスタンス属性: 特になし。
              \item 明示的に送出する例外: 明示的raiseはない。
              \item 制御構造の規模: 条件分岐 2、ループ 0、try 0
            \end{itemize}
            \paragraph{この定義を読むためのPython構文}
            \begin{itemize}[leftmargin=1.6em]
            \item 第1引数selfは、このメソッドを呼んだインスタンス自身を受け取る慣習名である。
\item with文はcontext managerに開始・終了処理を任せ、ファイルなどの資源を確実に解放する。
            \end{itemize}
            \paragraph{上から順の処理}
            \begin{enumerate}[leftmargin=1.8em]
            \item data に list(msg.data) の結果を代入する。
\item 条件 len(data) < 8 を判定し、真なら内部処理を行う。
\item    を返す。
\item with 文で self.\_lock を管理しながら処理する。
\item   self.\_state['batt\_voltage\_v'] に float(data[0]) の結果を代入する。
\item   self.\_state['batt\_current\_a'] に float(data[1]) の結果を代入する。
\item   self.\_set\_estimated\_soc(...) を実行する。
\item   self.\_state['motor\_w'] に float(data[3]) の結果を代入する。
\item   self.\_state['motor\_a'] に float(data[4]) の結果を代入する。
\item   self.\_state['solar\_w'] に float(data[5]) の結果を代入する。
\item   self.\_state['model\_speed\_kmh'] に float(data[6]) の結果を代入する。
\item   self.\_state['wheel\_w'] に float(data[7]) の結果を代入する。
\item   条件 len(data) >= 9 を判定し、真なら内部処理を行う。
\item     self.\_state['pack\_w'] に float(data[8]) の結果を代入する。
\item   self.\_state['ts'] に time.time() の結果を代入する。
            \end{enumerate}
            \paragraph{代表コード断片}
            \begin{lstlisting}[language=Python]
    def _on_metrics(self, msg: Float32MultiArray):
        data = list(msg.data)
        if len(data) < 8:
            return
        with self._lock:
            self._state['batt_voltage_v'] = float(data[0])
            self._state['batt_current_a'] = float(data[1])
            self._set_estimated_soc(data[2])
            self._state['motor_w'] = float(data[3])
            self._state['motor_a'] = float(data[4])
            self._state['solar_w'] = float(data[5])
            # Metrics carry the speed used for model evaluation, not the lower
            # controller command. Keep it separate so callback ordering cannot
            # overwrite /planner/speed_cmd on the dashboard.
            self._state['model_speed_kmh'] = float(data[6])
            self._state['wheel_w'] = float(data[7])
            if len(data) >= 9:
                self._state['pack_w'] = float(data[8])
            self._state['ts'] = time.time()
\end{lstlisting}

\subsection{L333 関数 DashboardNode.\_on\_upper\_plan}
            \begin{itemize}[leftmargin=1.6em]
              \item 行範囲: L333--L342
              \item 所属: \path|DashboardNode|
              \item 定義: \path|_on_upper_plan(self, msg: Float32MultiArray)|
              \item docstring: 明示されていない。
              \item このブロックが直接呼ぶ主な関数・メソッド: float, len, list, time
              \item 戻り値の要点: 戻り値が無いか、return 文が明示されていない。
              \item 主なローカル代入名: data, dt, speeds, v
              \item 読み取る主なインスタンス属性: self.\_lock, self.\_state
              \item 更新する主なインスタンス属性: 特になし。
              \item 明示的に送出する例外: 明示的raiseはない。
              \item 制御構造の規模: 条件分岐 1、ループ 0、try 0
            \end{itemize}
            \paragraph{この定義を読むためのPython構文}
            \begin{itemize}[leftmargin=1.6em]
            \item 第1引数selfは、このメソッドを呼んだインスタンス自身を受け取る慣習名である。
\item 内包表記は、反復・条件・値生成を一つの式で記述してcollectionを作る。
\item with文はcontext managerに開始・終了処理を任せ、ファイルなどの資源を確実に解放する。
            \end{itemize}
            \paragraph{上から順の処理}
            \begin{enumerate}[leftmargin=1.8em]
            \item data に list(msg.data) の結果を代入する。
\item 条件 len(data) < 2 を判定し、真なら内部処理を行う。
\item    を返す。
\item dt に float(data[0]) の結果を代入する。
\item speeds に [float(v) for v in data[1:]] の結果を代入する。
\item with 文で self.\_lock を管理しながら処理する。
\item   self.\_state['plan\_dt'] に dt の結果を代入する。
\item   self.\_state['plan\_upper'] に speeds の結果を代入する。
\item   self.\_state['ts'] に time.time() の結果を代入する。
            \end{enumerate}
            \paragraph{代表コード断片}
            \begin{lstlisting}[language=Python]
    def _on_upper_plan(self, msg: Float32MultiArray):
        data = list(msg.data)
        if len(data) < 2:
            return
        dt = float(data[0])
        speeds = [float(v) for v in data[1:]]
        with self._lock:
            self._state['plan_dt'] = dt
            self._state['plan_upper'] = speeds
            self._state['ts'] = time.time()
\end{lstlisting}

\subsection{L344 関数 DashboardNode.\_on\_lower\_plan}
            \begin{itemize}[leftmargin=1.6em]
              \item 行範囲: L344--L353
              \item 所属: \path|DashboardNode|
              \item 定義: \path|_on_lower_plan(self, msg: Float32MultiArray)|
              \item docstring: 明示されていない。
              \item このブロックが直接呼ぶ主な関数・メソッド: float, len, list, time
              \item 戻り値の要点: 戻り値が無いか、return 文が明示されていない。
              \item 主なローカル代入名: data, dt, speeds, v
              \item 読み取る主なインスタンス属性: self.\_lock, self.\_state
              \item 更新する主なインスタンス属性: 特になし。
              \item 明示的に送出する例外: 明示的raiseはない。
              \item 制御構造の規模: 条件分岐 1、ループ 0、try 0
            \end{itemize}
            \paragraph{この定義を読むためのPython構文}
            \begin{itemize}[leftmargin=1.6em]
            \item 第1引数selfは、このメソッドを呼んだインスタンス自身を受け取る慣習名である。
\item 内包表記は、反復・条件・値生成を一つの式で記述してcollectionを作る。
\item with文はcontext managerに開始・終了処理を任せ、ファイルなどの資源を確実に解放する。
            \end{itemize}
            \paragraph{上から順の処理}
            \begin{enumerate}[leftmargin=1.8em]
            \item data に list(msg.data) の結果を代入する。
\item 条件 len(data) < 2 を判定し、真なら内部処理を行う。
\item    を返す。
\item dt に float(data[0]) の結果を代入する。
\item speeds に [float(v) for v in data[1:]] の結果を代入する。
\item with 文で self.\_lock を管理しながら処理する。
\item   self.\_state['lower\_dt'] に dt の結果を代入する。
\item   self.\_state['plan\_lower'] に speeds の結果を代入する。
\item   self.\_state['ts'] に time.time() の結果を代入する。
            \end{enumerate}
            \paragraph{代表コード断片}
            \begin{lstlisting}[language=Python]
    def _on_lower_plan(self, msg: Float32MultiArray):
        data = list(msg.data)
        if len(data) < 2:
            return
        dt = float(data[0])
        speeds = [float(v) for v in data[1:]]
        with self._lock:
            self._state['lower_dt'] = dt
            self._state['plan_lower'] = speeds
            self._state['ts'] = time.time()
\end{lstlisting}

\subsection{L355 関数 DashboardNode.\_on\_sys\_state}
            \begin{itemize}[leftmargin=1.6em]
              \item 行範囲: L355--L356
              \item 所属: \path|DashboardNode|
              \item 定義: \path|_on_sys_state(self, msg: String)|
              \item docstring: 明示されていない。
              \item このブロックが直接呼ぶ主な関数・メソッド: \_update, str
              \item 戻り値の要点: 戻り値が無いか、return 文が明示されていない。
              \item 主なローカル代入名: 特になし。
              \item 読み取る主なインスタンス属性: self.\_update
              \item 更新する主なインスタンス属性: 特になし。
              \item 明示的に送出する例外: 明示的raiseはない。
              \item 制御構造の規模: 条件分岐 0、ループ 0、try 0
            \end{itemize}
            \paragraph{この定義を読むためのPython構文}
            \begin{itemize}[leftmargin=1.6em]
            \item 第1引数selfは、このメソッドを呼んだインスタンス自身を受け取る慣習名である。
            \end{itemize}
            \paragraph{上から順の処理}
            \begin{enumerate}[leftmargin=1.8em]
            \item self.\_update(...) を実行する。
            \end{enumerate}
            \paragraph{代表コード断片}
            \begin{lstlisting}[language=Python]
    def _on_sys_state(self, msg: String):
        self._update('system_state', str(msg.data))
\end{lstlisting}

\subsection{L358 関数 DashboardNode.\_on\_sys\_diag}
            \begin{itemize}[leftmargin=1.6em]
              \item 行範囲: L358--L359
              \item 所属: \path|DashboardNode|
              \item 定義: \path|_on_sys_diag(self, msg: String)|
              \item docstring: 明示されていない。
              \item このブロックが直接呼ぶ主な関数・メソッド: \_update, str
              \item 戻り値の要点: 戻り値が無いか、return 文が明示されていない。
              \item 主なローカル代入名: 特になし。
              \item 読み取る主なインスタンス属性: self.\_update
              \item 更新する主なインスタンス属性: 特になし。
              \item 明示的に送出する例外: 明示的raiseはない。
              \item 制御構造の規模: 条件分岐 0、ループ 0、try 0
            \end{itemize}
            \paragraph{この定義を読むためのPython構文}
            \begin{itemize}[leftmargin=1.6em]
            \item 第1引数selfは、このメソッドを呼んだインスタンス自身を受け取る慣習名である。
            \end{itemize}
            \paragraph{上から順の処理}
            \begin{enumerate}[leftmargin=1.8em]
            \item self.\_update(...) を実行する。
            \end{enumerate}
            \paragraph{代表コード断片}
            \begin{lstlisting}[language=Python]
    def _on_sys_diag(self, msg: String):
        self._update('system_diag', str(msg.data))
\end{lstlisting}

\subsection{L361 関数 DashboardNode.\_on\_mpc\_state}
            \begin{itemize}[leftmargin=1.6em]
              \item 行範囲: L361--L362
              \item 所属: \path|DashboardNode|
              \item 定義: \path|_on_mpc_state(self, msg: String)|
              \item docstring: 明示されていない。
              \item このブロックが直接呼ぶ主な関数・メソッド: \_update, str
              \item 戻り値の要点: 戻り値が無いか、return 文が明示されていない。
              \item 主なローカル代入名: 特になし。
              \item 読み取る主なインスタンス属性: self.\_update
              \item 更新する主なインスタンス属性: 特になし。
              \item 明示的に送出する例外: 明示的raiseはない。
              \item 制御構造の規模: 条件分岐 0、ループ 0、try 0
            \end{itemize}
            \paragraph{この定義を読むためのPython構文}
            \begin{itemize}[leftmargin=1.6em]
            \item 第1引数selfは、このメソッドを呼んだインスタンス自身を受け取る慣習名である。
            \end{itemize}
            \paragraph{上から順の処理}
            \begin{enumerate}[leftmargin=1.8em]
            \item self.\_update(...) を実行する。
            \end{enumerate}
            \paragraph{代表コード断片}
            \begin{lstlisting}[language=Python]
    def _on_mpc_state(self, msg: String):
        self._update('mpc_state', str(msg.data))
\end{lstlisting}

\subsection{L364 関数 DashboardNode.\_on\_health}
            \begin{itemize}[leftmargin=1.6em]
              \item 行範囲: L364--L365
              \item 所属: \path|DashboardNode|
              \item 定義: \path|_on_health(self, msg: Float32)|
              \item docstring: 明示されていない。
              \item このブロックが直接呼ぶ主な関数・メソッド: \_update, float
              \item 戻り値の要点: 戻り値が無いか、return 文が明示されていない。
              \item 主なローカル代入名: 特になし。
              \item 読み取る主なインスタンス属性: self.\_update
              \item 更新する主なインスタンス属性: 特になし。
              \item 明示的に送出する例外: 明示的raiseはない。
              \item 制御構造の規模: 条件分岐 0、ループ 0、try 0
            \end{itemize}
            \paragraph{この定義を読むためのPython構文}
            \begin{itemize}[leftmargin=1.6em]
            \item 第1引数selfは、このメソッドを呼んだインスタンス自身を受け取る慣習名である。
            \end{itemize}
            \paragraph{上から順の処理}
            \begin{enumerate}[leftmargin=1.8em]
            \item self.\_update(...) を実行する。
            \end{enumerate}
            \paragraph{代表コード断片}
            \begin{lstlisting}[language=Python]
    def _on_health(self, msg: Float32):
        self._update('system_health', float(msg.data))
\end{lstlisting}

\subsection{L367 関数 DashboardNode.\_on\_gps}
            \begin{itemize}[leftmargin=1.6em]
              \item 行範囲: L367--L371
              \item 所属: \path|DashboardNode|
              \item 定義: \path|_on_gps(self, msg: NavSatFix)|
              \item docstring: 明示されていない。
              \item このブロックが直接呼ぶ主な関数・メソッド: float, time
              \item 戻り値の要点: 戻り値が無いか、return 文が明示されていない。
              \item 主なローカル代入名: 特になし。
              \item 読み取る主なインスタンス属性: self.\_lock, self.\_state
              \item 更新する主なインスタンス属性: 特になし。
              \item 明示的に送出する例外: 明示的raiseはない。
              \item 制御構造の規模: 条件分岐 0、ループ 0、try 0
            \end{itemize}
            \paragraph{この定義を読むためのPython構文}
            \begin{itemize}[leftmargin=1.6em]
            \item 第1引数selfは、このメソッドを呼んだインスタンス自身を受け取る慣習名である。
\item with文はcontext managerに開始・終了処理を任せ、ファイルなどの資源を確実に解放する。
            \end{itemize}
            \paragraph{上から順の処理}
            \begin{enumerate}[leftmargin=1.8em]
            \item with 文で self.\_lock を管理しながら処理する。
\item   self.\_state['gps\_lat'] に float(msg.latitude) の結果を代入する。
\item   self.\_state['gps\_lon'] に float(msg.longitude) の結果を代入する。
\item   self.\_state['ts'] に time.time() の結果を代入する。
            \end{enumerate}
            \paragraph{代表コード断片}
            \begin{lstlisting}[language=Python]
    def _on_gps(self, msg: NavSatFix):
        with self._lock:
            self._state['gps_lat'] = float(msg.latitude)
            self._state['gps_lon'] = float(msg.longitude)
            self._state['ts'] = time.time()
\end{lstlisting}

\subsection{L373 関数 DashboardNode.\_on\_path}
            \begin{itemize}[leftmargin=1.6em]
              \item 行範囲: L373--L377
              \item 所属: \path|DashboardNode|
              \item 定義: \path|_on_path(self, msg: Path)|
              \item docstring: 明示されていない。
              \item このブロックが直接呼ぶ主な関数・メソッド: len, time
              \item 戻り値の要点: 戻り値が無いか、return 文が明示されていない。
              \item 主なローカル代入名: 特になし。
              \item 読み取る主なインスタンス属性: self.\_lock, self.\_state
              \item 更新する主なインスタンス属性: 特になし。
              \item 明示的に送出する例外: 明示的raiseはない。
              \item 制御構造の規模: 条件分岐 0、ループ 0、try 0
            \end{itemize}
            \paragraph{この定義を読むためのPython構文}
            \begin{itemize}[leftmargin=1.6em]
            \item 第1引数selfは、このメソッドを呼んだインスタンス自身を受け取る慣習名である。
\item with文はcontext managerに開始・終了処理を任せ、ファイルなどの資源を確実に解放する。
            \end{itemize}
            \paragraph{上から順の処理}
            \begin{enumerate}[leftmargin=1.8em]
            \item with 文で self.\_lock を管理しながら処理する。
\item   self.\_state['path\_points'] に len(msg.poses) の結果を代入する。
\item   self.\_state['ts'] に time.time() の結果を代入する。
            \end{enumerate}
            \paragraph{代表コード断片}
            \begin{lstlisting}[language=Python]
    def _on_path(self, msg: Path):
        # Keep last path length as a simple sanity signal
        with self._lock:
            self._state['path_points'] = len(msg.poses)
            self._state['ts'] = time.time()
\end{lstlisting}

\subsection{L379 関数 DashboardNode.\_init\_dummy}
            \begin{itemize}[leftmargin=1.6em]
              \item 行範囲: L379--L393
              \item 所属: \path|DashboardNode|
              \item 定義: \path|_init_dummy(self, path: str, rate_hz: float)|
              \item docstring: 明示されていない。
              \item このブロックが直接呼ぶ主な関数・メソッド: create\_timer, error, get\_logger, info, len, max, read\_csv, to\_dict
              \item 戻り値の要点: 戻り値が無いか、return 文が明示されていない。
              \item 主なローカル代入名: df, period
              \item 読み取る主なインスタンス属性: self.\_dummy\_rows, self.\_tick\_dummy, self.create\_timer, self.get\_logger
              \item 更新する主なインスタンス属性: self.\_dummy\_idx, self.\_dummy\_rows, self.\_dummy\_timer
              \item 明示的に送出する例外: 明示的raiseはない。
              \item 制御構造の規模: 条件分岐 1、ループ 0、try 1
            \end{itemize}
            \paragraph{この定義を読むためのPython構文}
            \begin{itemize}[leftmargin=1.6em]
            \item 第1引数selfは、このメソッドを呼んだインスタンス自身を受け取る慣習名である。
\item try文は例外が起き得る処理と、異常時または終了時の経路を分ける。
            \end{itemize}
            \paragraph{上から順の処理}
            \begin{enumerate}[leftmargin=1.8em]
            \item 例外処理を伴う try ブロックを実行する。
\item   Import 文を実行する。
\item   df に pd.read\_csv(path) の結果を代入する。
\item   Exceptionを捕捉した場合:
\item   self.get\_logger().error(...) を実行する。
\item    を返す。
\item 条件 df.empty を判定し、真なら内部処理を行う。
\item   self.get\_logger().error(...) を実行する。
\item    を返す。
\item self.\_dummy\_rows に df.to\_dict(orient='records') の結果を代入する。
\item self.\_dummy\_idx に 0 の結果を代入する。
\item period に 1.0 / max(rate\_hz, 0.5) の結果を代入する。
\item self.\_dummy\_timer に self.create\_timer(period, self.\_tick\_dummy) の結果を代入する。
\item self.get\_logger().info(...) を実行する。
            \end{enumerate}
            \paragraph{代表コード断片}
            \begin{lstlisting}[language=Python]
    def _init_dummy(self, path: str, rate_hz: float):
        try:
            import pandas as pd
            df = pd.read_csv(path)
        except Exception as exc:
            self.get_logger().error(f'Failed to load dummy_csv: {exc}')
            return
        if df.empty:
            self.get_logger().error('dummy_csv is empty.')
            return
        self._dummy_rows = df.to_dict(orient='records')
        self._dummy_idx = 0
        period = 1.0 / max(rate_hz, 0.5)
        self._dummy_timer = self.create_timer(period, self._tick_dummy)
        self.get_logger().info(f'Dummy mode enabled: {path} ({len(self._dummy_rows)} rows)')
\end{lstlisting}

\subsection{L395 関数 DashboardNode.\_tick\_dummy}
            \begin{itemize}[leftmargin=1.6em]
              \item 行範囲: L395--L406
              \item 所属: \path|DashboardNode|
              \item 定義: \path|_tick_dummy(self)|
              \item docstring: 明示されていない。
              \item このブロックが直接呼ぶ主な関数・メソッド: float, hasattr, items, len, time
              \item 戻り値の要点: 戻り値が無いか、return 文が明示されていない。
              \item 主なローカル代入名: key, row, val
              \item 読み取る主なインスタンス属性: self.\_dummy\_idx, self.\_dummy\_rows, self.\_lock, self.\_state
              \item 更新する主なインスタンス属性: self.\_dummy\_idx
              \item 明示的に送出する例外: 明示的raiseはない。
              \item 制御構造の規模: 条件分岐 1、ループ 1、try 1
            \end{itemize}
            \paragraph{この定義を読むためのPython構文}
            \begin{itemize}[leftmargin=1.6em]
            \item 第1引数selfは、このメソッドを呼んだインスタンス自身を受け取る慣習名である。
\item with文はcontext managerに開始・終了処理を任せ、ファイルなどの資源を確実に解放する。
\item try文は例外が起き得る処理と、異常時または終了時の経路を分ける。
            \end{itemize}
            \paragraph{上から順の処理}
            \begin{enumerate}[leftmargin=1.8em]
            \item 条件 not hasattr(self, '\_dummy\_rows') or not self.\_dummy\_rows を判定し、真なら内部処理を行う。
\item    を返す。
\item row に self.\_dummy\_rows[self.\_dummy\_idx] の結果を代入する。
\item self.\_dummy\_idx に (self.\_dummy\_idx + 1) \% len(self.\_dummy\_rows) の結果を代入する。
\item with 文で self.\_lock を管理しながら処理する。
\item   row.items() を順に走査し、各要素を (key, val) に入れて処理する。
\item     例外処理を伴う try ブロックを実行する。
\item       self.\_state[key] に float(val) の結果を代入する。
\item       Exceptionを捕捉した場合:
\item       self.\_state[key] に val の結果を代入する。
\item   self.\_state['ts'] に time.time() の結果を代入する。
            \end{enumerate}
            \paragraph{代表コード断片}
            \begin{lstlisting}[language=Python]
    def _tick_dummy(self):
        if not hasattr(self, '_dummy_rows') or not self._dummy_rows:
            return
        row = self._dummy_rows[self._dummy_idx]
        self._dummy_idx = (self._dummy_idx + 1) % len(self._dummy_rows)
        with self._lock:
            for key, val in row.items():
                try:
                    self._state[key] = float(val)
                except Exception:
                    self._state[key] = val
            self._state['ts'] = time.time()
\end{lstlisting}

\subsection{L409 関数 main}
            \begin{itemize}[leftmargin=1.6em]
              \item 行範囲: L409--L414
              \item 所属: module直下
              \item 定義: \path|main()|
              \item docstring: 明示されていない。
              \item このブロックが直接呼ぶ主な関数・メソッド: DashboardNode, destroy\_node, init, shutdown, spin
              \item 戻り値の要点: 戻り値が無いか、return 文が明示されていない。
              \item 主なローカル代入名: node
              \item 読み取る主なインスタンス属性: 特になし。
              \item 更新する主なインスタンス属性: 特になし。
              \item 明示的に送出する例外: 明示的raiseはない。
              \item 制御構造の規模: 条件分岐 0、ループ 0、try 0
            \end{itemize}
            \paragraph{この定義を読むためのPython構文}
            \begin{itemize}[leftmargin=1.6em]
            \item 特別な構文上の補足は少ない。
            \end{itemize}
            \paragraph{上から順の処理}
            \begin{enumerate}[leftmargin=1.8em]
            \item rclpy.init(...) を実行する。
\item node に DashboardNode() の結果を代入する。
\item rclpy.spin(...) を実行する。
\item node.destroy\_node(...) を実行する。
\item rclpy.shutdown(...) を実行する。
            \end{enumerate}
            \paragraph{代表コード断片}
            \begin{lstlisting}[language=Python]
def main():
    rclpy.init()
    node = DashboardNode()
    rclpy.spin(node)
    node.destroy_node()
    rclpy.shutdown()
\end{lstlisting}

        \subsection{設定値と入力パラメータ}
            \begin{longtable}{p{0.07\linewidth} p{0.35\linewidth} p{0.45\linewidth}}
            \toprule
            行 & 名前 & 既定値・補足 \\
            \midrule
            58 & \path|host| & 0.0.0.0 \\
59 & \path|port| & 8080 \\
60 & \path|static_dir| &  \\
61 & \path|dummy_csv| &  \\
62 & \path|dummy_rate_hz| & 5.0 \\
            \bottomrule
            \end{longtable}

        \subsection{ROS topic I/O}
            \begin{longtable}{p{0.07\linewidth} p{0.18\linewidth} p{0.34\linewidth} p{0.28\linewidth}}
            \toprule
            行 & 種別 & topic & callback / 補足 \\
            \midrule
            117 & subscription & \path|/planner/speed_cmd| & self.\_on\_speed\_cmd \\
118 & subscription & \path|/planner/upper_speed_cmd| & self.\_on\_upper\_speed\_cmd \\
119 & subscription & \path|/vehicle/speed_kmh| & self.\_on\_speed\_meas \\
120 & subscription & \path|/planner/throttle_cmd_pct| & self.\_on\_throttle\_cmd \\
121 & subscription & \path|/planner/drive_mode| & self.\_on\_drive\_mode \\
122 & subscription & \path|/vehicle/batt_soc| & self.\_on\_soc \\
123 & subscription & \path|/vehicle/batt_temp_c| & self.\_on\_tb \\
124 & subscription & \path|/vehicle/batt_current_a| & self.\_on\_ibatt \\
125 & subscription & \path|/vehicle/batt_voltage_v| & self.\_on\_vbatt \\
126 & subscription & \path|/vehicle/s_km| & self.\_on\_s\_km \\
127 & subscription & \path|/planner/status| & self.\_on\_status \\
128 & subscription & \path|/planner/env| & self.\_on\_env \\
129 & subscription & \path|/planner/metrics| & self.\_on\_metrics \\
130 & subscription & \path|/planner/upper_plan| & self.\_on\_upper\_plan \\
131 & subscription & \path|/planner/lower_plan| & self.\_on\_lower\_plan \\
132 & subscription & \path|/system/state| & self.\_on\_sys\_state \\
133 & subscription & \path|/system/diag| & self.\_on\_sys\_diag \\
134 & subscription & \path|/system/mpc_state| & self.\_on\_mpc\_state \\
135 & subscription & \path|/system/health| & self.\_on\_health \\
136 & subscription & \path|/sim/gps| & self.\_on\_gps \\
137 & subscription & \path|/planner/trajectory| & self.\_on\_path \\
            \bottomrule
            \end{longtable}







        \section{処理の流れ}
        \begin{enumerate}[leftmargin=1.7em]
        \item 初期化時に設定値や入力パスを読み込む。
\item publisher / subscription / timer を準備する。
\item timer callback 周期で主処理を進める。
        \end{enumerate}

        \section{このファイルの位置づけ}
        この資料群では、\path|mpc_solarcar/dashboard_node.py| を
        \textbf{runtime node}の中核として扱う。
        したがって、ここで扱う処理は単独で完結せず、
        前段の profile / launch / telemetry と後段の planner / logger / report へ繋がる。

        \end{document}
